\documentclass[a4paper,12pt]{report}
\setcounter{secnumdepth}{5}
\setcounter{tocdepth}{3}
\input{LaTeX-ToolKit/template.tex}
\begin{document}
\title{Ordinary Differential Equation}
\author{沈威宇}
\date{\temtoday}
\titletocdoc
\ch{Ordinary Differential Equation (常微分方程) and Dynamical System (動態系統 or 動力系統)}
\section{Ordinary Differential Equation}
\ssc{Definition}
\sssc{Ordinary Differential Equation (ODE) and System of ODEs (常微分方程組)}
An $n$th-order ODE or system of ODEs of a unknown function or dependent variable unknown functions or dependent variables $\mb{y}\colon\bbR\to\bbR^m$ with respect to an independent variable $x\in\bbR$ can be written in the form, called normal form,
\[\mb{y}^{(n)}=F(x,mb{y},mb{y}',mb{y}'',\ldots,\mb{y}^{(n-1)}),\]
in which $F\colon\Omega\subseteq\bbR\times\bbR^{m\times n}\to\bbR^m$ is a function.
\sssc{Initial-value problem (IVP) (初值問題)}
An $n$th-order initial-value problem is to solve an $n$th-order ODE or system of ODEs
\[\mb{y}^{(n)}=F(x,mb{y},mb{y}',mb{y}'',\ldots,\mb{y}^{(n-1)})\]
that subjects to an condition called initial condition (IC) (初始條件) or conditions called initial conditions at a point $x_0\in\bbR$ that can be written in the form
\[\mb{y}^{(i)}(x_0)=y_i,\quad\forall i\in\bbN_{0,<n}\]
and such that $(t,y_0,\ldots,y_{n-1})\in\Omega$.

A function $\mb{y}\colon\bbR\to\bbR^m$ that complies with the system of ODEs and the ICs is called a solution or solution vector.

The interval of existence of an IVP is the largest open interval containing $x_0$ where a solution of the IVP exists.

The interval of uniqueness of an IVP is the largest open interval containing $x_0$ where a unique solution of the IVP exists.
\sssc{Boundary-value problem (BVP) (邊值問題)}
An $n$th-order initial-value problem is to solve an $n$th-order ODE or system of ODEs that subjects to a condition called boundary condition (BC) (邊界條件) or conditions called boundary conditions.

A function $\mb{y}\colon\bbR\to\bbR^m$ that complies with the system of ODEs and the BCs is called a solution or solution vector.
\sssc{Well-posed problem}
A well-posed problem is a problem whose local solution satisfies the following properties:
\bit
\item It exists;
\item It is unique;
\item Its behavior changes continuously with the auxiliary conditions, such as initial or boundary values.
\eit
\ssc{Solutions of ODEs}
\sssc{Solution or flow of an ODE}
Any function $f$ that is defined on an interval $I$, called interval of definition, the interval of existence, the interval of validity, or the domain of the solution, and of class $C^n$ on $I$, and when substituted into an $n$th-order ordinary differential equation reduces the equation to an identity, is said to be a solution (aka flow) of the equation on $I$.

An ODE does not necessarily have to possess a solution.
\sssc{Solution curve or trajectory}
The graph of a solution $f$ on its interval of definition of an ODE is called a solution curve or a trajectory.
\sssc{Explicit solution}
An explicit solution of an ODE of a dependent variable $y$ with respect to an independent variable $x$ is in the form $y=f(x)$, where $f(x)$ is a function of $x$.
\sssc{Implicit solution}
A relation $G(x, y) = 0$ is said to be an implicit solution of an ODE on an interval $I$ if there exists at least one function $f$ that satisfes the relation $G(x, y) = 0$ as well as the ODE on $I$.
\sssc{Parametric family or parameterized family}
A parametric family or a parameterized family is a family of objects whose differences depend only on the chosen values for a set of parameters.
\sssc{Degrees of freedom}
The degrees of freedom of a system is the number of parameters of the system that may vary independently.
\sssc{Parametric family of solutions}
An $n$-parameter family of solutions of an ODE of a dependent variable $y$ with respect to an independent variable $x$ is in the form $y=f(x,c_1,c_2,\dots c_n)$ (explicit) or $G(x,y,c_1,c_2,\dots c_n)=0$ (implicit), in which $c_1,c_2,\dots c_n$ are parameters that are arbitrary given that the solution obtained is a solution of the ODE. The arbitrary parameters are usually denoted as $A,B,\ldots$.
\sssc{General solution (通解)}
If every solution of an ODE on an interval $I$ can be obtained from an $n$-parameter family of solutions by appropriate choices of the parameters, we then say that that family of solutions is the general solution of the ODE.
\sssc{Particular solution (特解)}
A particular solution of an ODE that is free of parameters is called a particular solution.

If a system given by an ODE have $n$ degrees of freedom, one may obtain a determined particular solution by imposing $n-1$ non contradictory constraints, leaving one degree of freedom determined by the ODE.
\sssc{Singular solution (奇異解)}
Sometimes a differential equation possesses a solution that is not a member of a family of solutions of the equation, that is, a solution that cannot be obtained by specializing any of the parameters in the family of solutions. Such an extra solution is called a singular solution. The independent variable value at which is called singular point.
\sssc{Solutions by substitutions or solutions by change of variables}
When solving an ODE, we often replace dependent variables with functions of the independent variable and new dependent variables where the new dependent variables are functions of the independent variable and old dependent variables to transform it into another ODE of the new dependent variables that is easier to solve.
\sssc{Solution of System of ODEs}
A solution of a system of ODEs involving the derivatives of $n$ unknown functions or dependent variables, $y_1,y_2,\dots y_n$ of a single independent variable $x$ is a $n$-tuple of sufficiently smooth functions of $x$ defined on a common interval of definition that satisfy all equations in the system.

A system of ODEs does not necessarily have to possess a solution.
\sssc{Orthogonal trajectory and isogonal trajectory}
An orthogonal trajectory of a family of curves is a curve that intersects each curve of the family orthogonally.

An isogonal trajectory of a family of curves is a curve that intersects each curve of the family at a same angle $\alpha$.
\ssc{Conversion of order to variable}
Any $n$th-order system of ODEs of $m$ variables $x_1,x_2,\ldots,x_m$ can be converted to a first-order system of ODEs of $m\cdot n$ variables by, for each variable $x$, assigning $(n-1)$ new variables defined as the first to $(n-1)$th derivatives of $x$.
\section{First-order ODE}
Unless otherwise specified, the variables and functions are in $\mathbb{R}$.
\ssc{Differential form}
The differential form of a first-order ODE $M(x,y)+N(x,y)\dv{y}{x}=0$ is
\[M(x,y)\dd{x}+N(x,y)\dd{y}=0.\]
\ssc{Slope}
\sssc{Slope function or rate function}
In a first-order ODE in the form $\dv{y}{x}=f(x,y)$, $f$ is called the slope function or rate function.
\sssc{Direction field or slope field (斜率場)}
In a first-order ODE in the form $\dv{y}{x}=f(x,y)$, if we evaluate $f$ over a rectangle grid of points and draw a lineal element at each point $(x,y)$ of the grid with slope $f(x,y)$, then the collection of all these lineal elements is called a direction field or a slope field of the ODE.
\sssc{Isocline}
In a first-order ODE in the form $\dv{y}{x}=f(x,y)$, the isocline for slope $m$ is defined as the family of points $\{(x,y)\mid f(x,y)=m\}$. In the method of isocline, we draw a lineal element at each point in the isocline for slope $m$ of $f$ with slope $m$.
\ssc{Autonomous System}
\sssc{Autonomous system (自治系統)}
A system of ODEs is called autonomous if it can be written in the form
\[\dv{y}{x}=f(y)\]
\sssc{Critical point (臨界點), equilibrium, equilibrium point (平衡點), or stationary point (駐點) of autonomous systems of ODEs}
For an autonomous system, we say that a point $c$ in the domain of $f$ is a critical point, equilibrium, equilibrium point, or stationary point iff $f(c)=0$.

If $c$ is a critical point, then $y=c$ is a constant solution, also called equilibrium solution, of the system. If also that $f$ is locally Lipschitz at $c$, then $y=c$ is the unique solution of the system through $c$.
TODO
\ssc{Orthogonal trajectories}
The family of the orthogonal trajectories of the family of solutions of the first-order ODE
\[M(x,y)\dd{x}+N(x,y)\dd{y}=0,\]
is the family of solutions of the first-order ODE
\[N(x,y)\dd{x}-M(x,y)\dd{y}=0.\]
\ssc{Existence and Uniqueness of Solution}
\sssc{Picard–Lindelöf theorem, Cauchy–Lipschitz theorem, or existence and uniqueness theorem (local version)}
Let $U$ be an open subset of $\bbR\times\bbR^n$ and $f\colon U\to\bbR^n$ be a function that is continuous in $t$ and locally Lipschitz in $y$, that is, for any $(t_0,y_0)\in U$, there exists a neighborhood $V\subseteq U$ of $(t_0,y_0)$ and a $L>0$ such that for all $(t,y_1),(t,y_2)\in V$, $\norm{f(t,y_1)-f(t,y_2)}\le L\norm{y_1-y_2}$.

For any $(t_0,y_0)\in U$, there exists some $\varepsilon>0$ such that the initial value problem
\[y'(t)=f(t,y(t)),\quad y(t_0)=y_0\]
has a unique solution $y(t)$ on the interval $[t_0-\varepsilon,t_0+\varepsilon]$.

Moreover, let $I_1,I_2\subseteq\bbR$ be two open intervals and $y_i\colon I_i\to\bbR^n$ for $i\in\{1,2\}$ two continuously differentiable functions such that $(t,y_i(t))\in U$ and
\[y_i'(t)=f(t,y_i(t))\]
for all $t\in I_i$, for $i\in\{1,2\}$. If $y_1(s)=y_2(s)$ for some $s\in I_1\cap I_2$, then $y_1(t)=y_2(t)$ for all $t\in I_1\cap I_2$.
\bpr
Let $D\subseteq U$ be a closed rectangle such that $(t_0,y_0)$ is in the interior of $D$. Since local Lipschitz on compact set implies global Lipschitz on it, $f$ restricted to $D$ is globally Lipschitz continuous in $y$, that is,
\[\exists L>0\colon\forall (t,y_1),(t,y_2)\in D\colon\norm{f(t,y_1)-f(t,y_2)}\le L\norm{y_1-y_2}.\]
Let $B_{\leq r}(a)$ denote closed ball of radius $r$ centered at $a$.

We first show that $f$ is continuous at $(t_0,y_0)$. Since $t\mapsto f(t,y)$ is continuous,
\[\forall (t_0,y_0)\in D,\varepsilon>0\colon\exists\delta>0\colon\norm{t-t_0}<\delta\implies\norm{f(t,y_0)-f(t_0,y_0)}<\frac{\varepsilon}{2}\]
For any $t\in B_{\leq\delta}(t_0)$, $y\in B_{\le\frac{\varepsilon}{2L}}(y_0)$,
\[\norm{f(t,y)-f(t_0,y_0)}\le\norm{f(t,y)-f(t,y_0)}+\norm{f(t,y_0)-f(t_0,y_0)}<\varepsilon.\]
Fix any such $L>0$ and any $b>0$ such that
\[C_o\coloneq B_{\leq\frac{1}{2L}}(t_0)\times B_{\le b}(y_0)\subseteq D.\]
Such set exists because $(t_0,y_0)$ is in the interior of $D$.

Since $C_o$ is compact, there exists a real number
\[M_o=\sup_{(t,y)\in C_o}\norm{f(t,y)}.\]
Define
\[a\coloneq\frac{1}{2}\min\qty{\frac{1}{L},\frac{b}{M_o}}\]
and define
\[I\coloneq B_{\leq a}(t_0)\]
\[B\coloneq B_{\leq b}(y_0)\]
\[C\coloneq I\times B\subseteq C_o\]
\[M=\sup_{(t,y)\in C}\norm{f(t,y)}\le M_o\]
Consider the function space $\calC(I,B)$ of continuous (and thus bounded since $I$ is closed) function from $I$ to $B$ equipped with uniform norm $\|\cdot\|_{\infty}$ and metric induced by it. We want to prove that $\calC(I,B)$ is complete.

Assume $(y_n)_{n\in\bbN}$ is a Cauchy sequence in $\calC(I,B)$.
\[\forall\varepsilon>0\colon\exists N\in\bbN\colon m,n\ge N\implies\norm{y_m-y_n}_{\infty}<\varepsilon\implies\sup_{t\in I}\norm{y_m(t)-y_n(t)}<\varepsilon\]
then
\[\exists y\in\calC(I,\bbR)\colon\forall\varepsilon>0\colon\exists N\in\bbN\colon n\ge N\implies\sup_{t\in I}\norm{y_n(t)-y(t)}<\varepsilon,\]
that is, $(y_n)$ converges uniformly to some $y\in\calC(I,\bbR)$. Since $B$ is closed, $y(t)\in B$. Hence $y\in\calC(I,B)$.

Integrate both sides of the original equation
\[y(t)=y_0+\int_{t_0}^tf(s,y(s))\dd{s}.\]
Define Picard operator $\Gamma$
\[(\Gamma y)(t)=y_0+\int_{t_0}^tf(s,y(s))\dd{s},\quad\forall y\in \calC(I,B),t\in I.\]
Then, the statement is equivalent to that there exists a unique fixed point of $\Gamma$, which is the solution $y(t)$ of the original equation. To apply the Banach fixed-point theorem, we must show that $\Gamma$ is a contraction mapping over a non-empty complete metric space.

We first show that $\Gamma(\calC(I,B))\subseteq\calC(I,B)$, that is,
\[\|y-y_0\|_{\infty}\le b\implies\|\Gamma y-y_0\|\le b.\]
Let
\[t_m=\arg\max_{t\in I}\norm{\int_{t_0}^tf(s,y(s))\dd{s}}\]
\[\ba
  \|\Gamma y-y_0\|&=\norm{\int_{t_0}^tf(s,y(s))\dd{s}}_{\infty}\\
                  &=\norm{\int_{t_0}^{t_m}f(s,y(s))\dd{s}}\\
                  &\le\norm{\int_{t_0}^{t_m}\norm{f(s,y(s))}\dd{s}}\\
                  &\le\norm{\int_{t_0}^{t_m}M\dd{s}}\\
                  &\le M\norm{t_m-t_0}\\
                  &\le Ma\\
                  &\le b
\ea\]
Second, we want to show that
\[\exists q\in[0,1)\colon\forall y_1,y_2\in\calC(I,B)\colon\norm{\Gamma y_1-\Gamma y_2}_{\infty}\le q\norm{y_1-y_2}_{\infty}\]
We claim that it holds for $q=aL\le\frac{1}{2}$. Fix any $y_1,y_2\in\calC(I,B)$. Let
\[t_m=\arg\max_{t\in I}\norm{(\Gamma y_1-\Gamma y_2)(t)}\]
Then, using the definition of $\Gamma$,
\[\ba
  \norm{\Gamma y_1-\Gamma y_2}_{\infty}&=\norm{(\Gamma y_1-\Gamma y_2)(t_m)}\\
                                       &=\norm{\int_{t_0}^{t_m}f(s,y_1(s))-f(s,y_2(s))\dd{s}}\\
                                       &\le\norm{\int_{t_0}^{t_m}\norm{f(s,y_1(s))-f(s,y_2(s))}\dd{s}}\\
                                       &\le L\norm{\int_{t_0}^{t_m}\norm{y_1(s)-y_2(s)}\dd{s}}\\
                                       &\le L\norm{\int_{t_0}^{t_m}\norm{y_1-y_2}_{\infty}\dd{s}}\\
                                       &=L\norm{t_m-t_0}\norm{y_1-y_2}_{\infty}\\
                                       &\le aL\norm{y_1-y_2}_{\infty}\\
                                       &=q\norm{y_1-y_2}_{\infty}
\ea\]
Now for the second statement, we define
\[J\coloneq I_1\cap I_2\]
\[S\coloneq\{t\in J\colon y_1(t)=y_2(t)\}\]
Since $s\in S$, $S$ is non-empty. Since $y_1,y_2$ are continuous,
\[g(t)\coloneq y_1(t)-y_2(t)\]
is continuous. Therefore,
\[S=g^{-1}(\{\mb{0}\})\]
is closed in $J$.

Fix any $t_0\in S$ and write $y_0=y_1(t_0)=y_2(t_0)$. By the fist statement, there exists $\varepsilon>0$ such that the initial value problem
\[y'(t)=f(t,y(t)),\quad y(t_0)=y_0\]
has a unique solution $y(t)$ on the interval $[t_0-\varepsilon,t_0+\varepsilon]$, which implies that $[t_0-\varepsilon,t_0+\varepsilon]\cap J$ is a subset of $S$.

Assume for contradiction that $S\ne J$. Then, there exists an endpoint $x$ of $S$ such that a neighbourhood $X$ of $x$ must contains some point in $J\setminus S$. By picking $t_0=x$, there exists $\varepsilon>0$ such that $[t_0-\varepsilon,t_0+\varepsilon]\cap J$ is a subset of $S$, but it must also contains some point in $J\setminus S$, which is contradictory.
\epr
\sssc{Peano existence theorem (皮亞諾存在性定理), Peano theorem (皮亞諾定理), or Cauchy–Peano theorem (柯西-皮亞諾定理)}
TODO
Let $D$ be an open subset of $\mathbb{R}\times\mathbb{R}^n$ and $f\colon D\to \mathbb {R}$ be a continuous function, then every initial value problem given by explicit first-order ODE $y'(t)=f\left(t,y(t)\right)$ defined on $D$ and initial condition $y\left(t_{0}\right)=y_{0}$ with $(t_{0},y_{0})\in D$ has a local solution $z\colon I\to \mathbb {R}^n$ where $I$ is a neighbourhood of $t_0$ in $\mathbb {R} $.
\sssc{Carathéodory's existence theorem}
TODO
Let
\[D=\{(t,y)\mid|t-T|\leq a\land\forall\tx{\ integer\ }1\leq i\leq n\colon |y_i-Y_i|\leq b_i\}\]
be a rectangle in $\mathbb{R}\times\mathbb{R}^n$ where $t,T\in\mathbb{R}$, $y=(y_1, y_2,\dots  y_n)\in\mathbb{R}^n$, $Y=(Y_1, Y_2,\dots  Y_n)\in\mathbb{R}^n$, and constant $a,b_i\in\mathbb{R}_{>0}$, and $f(t,y)\colon D\to \mathbb {R}$ be a function that is:
\begin{itemize}
\item continuous in $y$ for each fixed $t$,
\item Lebesgue-measurable in $t$ for each fixed $y$, and
\item such that there is a Lebesgue-integrable function $m\colon [T-a,T+a]\to [0,\infty )$ such that $|f(t,y)|\leq m(t)$ for all $(t,y)\in D$,
\end{itemize}
then every initial value problem given by explicit first-order ODE $y'(t)=f\left(t,y(t)\right)$ defined on $D$ and initial condition $y\left(t_{0}\right)=y_{0}$ with $(t_{0},y_{0})\in D$ has a local solution $z\colon I\to \mathbb {R}^n$ where $I$ is a neighbourhood of $t_0$ in $\mathbb{R}$.

Let $I$ be an open interval of $\mathbb{R}$ and $f(t,y)\colon I\times\mathbb{R}^n\to\mathbb{R}$ where $t\in\mathbb{R}$, $y=(y_1, y_2,\dots  y_n)\in\mathbb{R}^n$ be a function that is:
\begin{itemize}
\item continuous in $y$ for each fixed $t$,
\item Lebesgue-measurable in $t$ for each fixed $y$, and
\item such that there is a Lebesgue-integrable function $m\colon I\to [0,\infty )$ such that $|f(t,y)|\leq m(t)$ for all $(t,y)\in I\times\mathbb{R}^n$,
\end{itemize}
then every initial value problem given by explicit first-order ODE $y'(t)=f\left(t,y(t)\right)$ defined on $I\times\mathbb{R}^n$ and initial condition $y\left(t_{0}\right)=y_{0}$ with $(t_{0},y_{0})\in I\times\mathbb{R}^n$ has a local solution $z\colon I\to \mathbb {R}^n$ where $I$ is a neighbourhood of $t_0$ in $\mathbb{R}$.
\sssc{Cauchy–Kovalevskaya theorem}
TODO
\ssc{Separable equation}
\sssc{Definition}
A first-order ODE in the form
\[\dv{y}{x}=g(x)h(y)\]
is said to be separable or have separable variables.
\sssc{Solutions by direct integration}
The solutions of
\[\dv{y}{x}=g(x)h(y)\]
are the family of solutions
\[\int\frac{1}{h(y)}\dd{y}=\int g(x)\dd{x}+C\]
and the constant solution
\[h(y)=0.\]
\sssc{Logistic differential equation}
IVP:
\[\dv{P}{t}=kP\qty(1-\frac{P}{M}),\quad P(0)=P_0.\]
Solutions:
\[P=\frac{M}{1+\frac{M-P_0}{P_0}e^{-kt}}.\]
\begin{proof}
\[\frac{1}{P\qty(1-\frac{P}{M})}\dd{P}=k\dd{t}\]
\[\int\frac{1}{P\qty(1-\frac{P}{M})}\dd{P}=\int k\dd{t}=kt+A\]
\[\frac{1}{P\qty(1-\frac{P}{M})}=\frac{M}{P(M-P)}=\frac{A}{P}+\frac{B}{M-P}\]
\[AM-AP+BP=M,\quad A=1,\quad B=1\]
\[\int\frac{1}{P\qty(1-\frac{P}{M})}\dd{P}=\int\frac{1}{P}+\frac{1}{M-P}\dd{P}=\ln\abs{P}-\ln\abs{M-P}+C\]
\[\ln\abs{P}-\ln\abs{M-P}=kt+A=\ln\abs{\frac{P}{M-P}}\]
\[\frac{P}{M-P}=Ae^{kt}\]
\[P=\frac{AMe^{kt}}{1+Ae^{kt}}\]
\[P_0=\frac{AM}{1+A}\]
\[A=\frac{P_0}{M-P_0}\]
\[P=\frac{\frac{P_0}{M-P_0}Me^{kt}}{1+\frac{P_0}{M-P_0}e^{kt}}=\frac{M}{1+\frac{M-P_0}{P_0}e^{-kt}}\]
\end{proof}
\sssc{Lotka–Volterra equations or predator–prey equations}
ODE:
\[\dv{x}{t}=x(\alpha-\beta y),\]
\[\dv{y}{t}=-y(\gamma-\delta x).\]
Implicit solutions:
\[\delta x-\gamma\ln(x)+\beta y-\alpha\ln(y)=C.\]
Equilibrium solutions:
\[x=0,\quad y=0,\]
and
\[x=\frac{\gamma}{\delta},\quad y=\frac{\alpha}{\beta}.\]
\begin{proof}
\[\dv{y}{t}=\dv{y}{x}\dv{x}{t}\]
\[\dv{y}{x}=-\frac{y}{x}\frac{\gamma-\delta x}{\alpha-\beta y}\]
\[\frac{\delta x-\gamma}{x}\dd{x}=-\frac{\beta y-\alpha}{y}\dd{y}\]
\[\delta x-\gamma\ln(x)+\beta y-\alpha\ln(y)=C\]
\end{proof}
\ssc{First-order linear equation}
\sssc{Definition}
A first-order ODE in the form
\[\dv{y}{x}+P(x)y=Q(x)\]
is said to be linear.
\sssc{Solutions by integrating factor}
Use the integrating factor:
\[\mu(x)=e^{\int P(x)\dd{x}}.\]
Multiply both sides with $\mu(x)$:
\[\mu(x)\dv{y}{x}+\mu(x)P(x)y=\mu(x)Q(x).\]
The left-hand side is exactly a derivative:
\[\dv{}{x}\qty(\mu(x)y(x))=\mu(x)Q(x).\]
Integrate:
\[\mu(x)y(x)=\int\mu(x)Q(x)\dd{x}+C.\]
\[y(x)=\frac{1}{\mu(x)}\qty(\int\mu(x)Q(x)\dd{x}+C).\]
\sssc{Remarks}
Occasionally, a first-order differential equation is not linear in one variable but is linear in the other variable. If a first-order ODE that can be wriiten in the form
\[\dv{y}{x}=\frac{1}{P(y)x+Q(y)}.\]
We can take its reciprocal as
\[\dv{x}{y}-P(y)x=Q(y),\]
which is recognized as linear in the variable $x$.
\ssc{Exact equation}
\sssc{Definition}
On a simply-connected open set, a first-order ODE in the form
\[M(x,y)\dd{x}+N(x,y)\dd{y}=0\]
where $\frac{\partial^2M}{\partial x\partial y}$ and $\frac{\partial^2N}{\partial x\partial y}$ is continuous and
\[\pdv{M}{y}=\pdv{N}{x}\]
is said to be exact, and its LHS is called exact differential.

By Schwarz's Theorem, the criterion holds iff there exists a function $F\colon R\to\bbR$ such that
\[\pdv{F}{x}=M,\quad\pdv{F}{y}=N,\]
and that $\frac{\partial^2F}{\partial x\partial y}$ is continuous. Such $F$ is called the potential function.
\sssc{Solutions by integration and then differentiation with respect to different variables}
First, integrate $M(x,y)$ with respect to $x$:
\[F(x,y)=\int M(x,y)\dd{x}+\phi(y),\]
where $\phi(y)$ is an arbitrary constant of integration that may depend on $y$.

Second, differentiate with respect to $y$:
\[\pdv{F}{y}(x,y)=N(x,y)=\pdv{}{y}\qty(\int M(x,y)\dd{x})+\phi'(y).\]

Third, integrate $\phi'(y)$ with respect to $y$ to find $\phi(y)$:
\[\phi(y)=\int\phi'(y)\dd{y}=\int N(x,y)-\pdv{}{y}\qty(\int M(x,y)\dd{x})\dd{y}.\]

Fourth, substitute $\phi(y)$ back to find the potential function $F(x,y)$. The implicit solutions is given by
\[F(x,y)=C,\]
where $C$ is an arbitrary constant.

Alternatively, first, integrate $N(x,y)$ with respect to $y$:
\[F(x,y)=\int N(x,y)\dd{y}+\phi(x),\]
where $\phi(x)$ is an arbitrary constant of integration that may depend on $x$.

Second, differentiate with respect to $x$:
\[\pdv{F}{x}(x,y)=M(x,y)=\pdv{}{x}\qty(\int N(x,y)\dd{y})+\phi'(x).\]

Third, integrate $\phi(x)$ with respect to $x$ to find $\phi(x)$:
\[\phi(x)=\int\phi'(x)\dd{x}=\int M(x,y)-\pdv{}{x}\qty(\int N(x,y)\dd{y})\dd{x}.\]

Fourth, substitute $\phi(x)$ back to find the potential function $F(x,y)$. The implicit solutions is given by
\[F(x,y)=C,\]
where $C$ is an arbitrary constant.
\ssc{Modified exact equation}
\sssc{General case}
On a simply-connected open set, for a first-order ODE in the form
\[M(x,y)\dd{x}+N(x,y)\dd{y}=0,\]
we define the integrating factor $\mu(x,y)$ such that
\[\mu(x,y)M(x,y)\dd{x}+\mu(x,y)N(x,y)\dd{y}=0\]
is exact and with potential function $F$ such that
\[\pdv{F}{x}=\mu(x,y)M(x,y),\quad\pdv{F}{y}=\mu(x,y)N(x,y),\]
that is,
\[\pdv{\mu M}{y}=\pdv{\mu N}{x}.\]
Solve the PDE to get $\mu$.
\sssc{Separable case}
If
\[\frac{\pdv{M}{y}-\pdv{N}{x}}{N}\]
depends only on $x$, we can construct integrating factor $\mu(x)$:
\[\mu(x)=e^{\int\frac{\pdv{M}{y}-\pdv{N}{x}}{N}\dd{x}}\]
such that
\[\pdv{\mu M}{y}=\pdv{\mu N}{x}=\mu\pdv{M}{y}.\]
If
\[\frac{\pdv{N}{x}-\pdv{M}{y}}{M}\]
depends only on $y$, we can construct integrating factor $\mu(y)$:
\[\mu(y)=e^{\int\frac{\pdv{N}{x}-\pdv{M}{y}}{M}\dd{x}}\]
such that
\[\pdv{\mu M}{y}=\pdv{\mu N}{x}=\mu\pdv{N}{x}.\]
\begin{proof}
\[\frac{\mu'(x)}{\mu(x)}=\frac{\pdv{M}{y}-\pdv{N}{x}}{N}.\]
This is a separable ODE of $\mu$ with respect to $x$.

Let
\[P(x)=\frac{\pdv{M}{y}-\pdv{N}{x}}{N}.\]
\[\mu^{-1}\dd{\mu}=P(x)\dd{x}.\]
\[\ln|\mu|=\int P(x)\dd{x}.\]
\[\mu=e^{\int P(x)\dd{x}}.\]

Similarly,
\[\frac{\mu'(y)}{\mu(y)}=\frac{\pdv{N}{x}-\pdv{M}{y}}{M}\]
for the second case.
\end{proof}
\ssc{Homogeneous equation}
\sssc{Definition}
A first-order ODE in the form
\[\dv{y}{x}=F(x,y)\]
with $F(x,y)$ be a homogeneous function of degree $0$, i.e.,
\[\forall t\ne0\colon F(tx,ty)=F(x,y),\]
is said to be homogeneous. 

A first-order ODE in the form
\[M(x,y)\dd{x}+N(x,y)\dd{y}=0,\]
where $M(x,y)$ and $N(x,y)$ are homogeneous functions of the same degree is homogeneous since $F(x,y)=-\frac{M(x,y)}{N(x,y)}$ is a homogeneous function of degree $0$.
\sssc{Solutions by substitutions}
For any homogeneous function $F(x,y)$ of $2$ variables, there exists a function $f$ of $1$ variable such that
\[\forall x\ne0,(x,y)\in D_F\colon F(x,y)=f\qty(\frac{y}{x})\]
Introduce the substitution
\[v=\frac{y}{x}\]
\[\dv{y}{x}=\dv{(xv)}{x}=v+x\dv{v}{x}=f(v),\]
which is a separable ODE.
\[\frac{\dd{v}}{f(v)-v}=\frac{\dd{x}}{x}.\]
\[\int\frac{\dd{v}}{f(v)-v}=\ln|x|+C.\]
Substituting $v=\frac{y}{x}$ back gives the solutions.

If $F$ is defined when $x=0$, we also have a implicit constant solution
\[F(0,y)=0.\]
\ssc{Bernoulli equation}
\sssc{Definition}
A first-order ODE in the form
\[\dv{y}{x}+P(x)y=Q(x)y^n,\quad n\in\mathbb{R}\setminus\{0,1\}\]
is called a Bernoulli equation.

When $n=0$ or $1$, it is a first-order linear equation.
\sssc{Solutions by substitutions}
Divide both sides by $y^n$:
\[y^{-n}\dv{y}{x}+P(x)y^{1-n}=Q(x).\]
Introduce the substitution
\[v=y^{1-n}\]
\[\dv{v}{x}=(1-n)y^{-n}\dv{y}{x}.\]
Substitute into the equation:
\[\frac{1}{1-n}\dv{v}{x}+P(x)v=Q(x).\]
\[\dv{v}{x}+(1-n)P(x)v=(1-n)Q(x),\]
which is linear.

Use the integrating factor:
\[\mu(x)=e^{\int (1-n)P(x)\dd{x}}.\]
Multiply both sides with $\mu(x)$:
\[\mu(x)\dv{v}{x}+\mu(x)(1-n)P(x)v=\mu(x)(1-n)Q(x).\]
The left-hand side is exactly a derivative:
\[\dv{}{x}\qty(\mu(x)v(x))=\mu(x)(1-n)Q(x).\]
Integrate:
\[\mu(x)v(x)=\int\mu(x)(1-n)Q(x)\dd{x}+C.\]
\[v(x)=\frac{1}{\mu(x)}\qty(\int\mu(x)(1-n)Q(x)\dd{x}+C).\]
Recall $v=y^{1-n}$, we obtain the family of general solutions.
\ssc{Riccati equation}
\sssc{Definition}
A first-order ODE in the form
\[\dv{y}{x}=q_0(x)+q_1(x)y+q_2(x)y^2,\]
where $q_2(x)$ is nonzero and differentiable, is called a Riccati equation.

If $q_2(x)=1$, it is called canonical form.
\sssc{Solutions by substitutions}
First convert to canonical form. Let
\[v=yq_2\]
\[\ba
v'&=y'q_2+yq_2'\\
  &=(q_0+q_1y+q_2y^2)q_2+v\frac{q_2'}{q_2}\\
  &=q_0q_2+v\qty(q_1+\frac{q_2'}{q_2})+v^2
\ea\]
Define
\[a(x)=q_0q_2\]
\[b(x)=q_1+\frac{q_2'}{q_2}\]
We obtain a canonical form Riccati equation:
\[\dv{v}{x}=a(x)+b(x)v+v^2\]
Second, convert it to a second-order homogeneous linear equation. Let
\[u=\exp(-\int v\dd{x})\]
\[u'=-uv\]
\[\ba
u''&=-u'v-v'u\\
   &=v^2u-(a+bv+v^2)u\\
   &=-au-bvu\\
   &=-au+bu'
\ea\]
We obtain a second-order homogeneous linear equation:
\[u''-bu'+au=0\]
Solve it to obtain $u$.
\[y=-q_2^{-1}\qty(\ln|u|)'\]
If $u_1,u_2$ are two linearly independent solution, the general solution of $y$ is
\[y=-q_2^{-1}\qty(\ln|c_1u_1+c_2u_2|)'\]
\ssc{Slope being function of linear function}
ODE:
\[\dv{y}{x}=f(u),\quad u=ax+by+c.\]
Solutions:
\[\int\frac{1}{a+bf(u)}\dd{u}=x+C.\]
\begin{proof}
\[\dv{u}{x}=a+b\dv{y}{x}=a+bf(u)\]
This is separable.
\[\int\frac{1}{a+bf(u)}\dd{u}=x+C\]
\end{proof}
\ssc{Slope being function of linear function over linear function}
ODE:
\[\dv{y}{x}=F\qty(\frac{ax+by+c}{px+qy+r}),\quad aq-bp\neq 0.\]
We can find a unique solution of $ax+by+c=px+qy+r=0$, $x=\alpha$ and $y=\beta$. Let
\[f(t)=F\qty(\frac{a+bt}{p+qt}),\]
and we can rewrite it as:
\[\dv{y}{x}=f(u),\quad u=\frac{y-\beta}{x-\alpha}.\]
\begin{proof}
\[y=u(x-\alpha)+\beta\]
\[ax+by+c=(a+bu)(x-\alpha)+a\alpha+b\beta+c=(a+bu)(x-\alpha)\]
\[px+qy+r=(p+qu)(x-\alpha)+p\alpha+q\beta+r=(p+qu)(x-\alpha)\]
\[\frac{ax+by+c}{px+qy+r}=\frac{a+bu}{p+qu}\]
\end{proof}
Solutions:
\[\int\frac{1}{f(u)-u}\dd{u}=\ln|x-\alpha|+C,\]
and possible constant or singular solutions:
\[f(u)=u.\]
\begin{proof}
\[y=(x-\alpha)u+\beta\]
\[\dv{y}{x}=\pdv{}{x}\qty((x-\alpha)u+\beta)+\pdv{}{u}\qty((x-\alpha)u+\beta)\dv{u}{x}=u+(x-\alpha)\dv{u}{x}\]
\[u+(x-\alpha)\dv{u}{x}=f(u)\]
\[\frac{1}{f(u)-u}\dd{u}=\frac{1}{x-\alpha}\dd{x}\]
This is separable.
\[\int\frac{1}{f(u)-u}\dd{u}=\ln|x-\alpha|+C\]
\end{proof}
\section{Linear ODE}
\ssc{Basics}
\sssc{Linear ODE}
An $n$th-order linear ODE of a unknown function or dependent variable $y\colon\bbR\to\bbR$ with respect to an independent variable $x\in\bbR$ can be written in the form, called normal form,
\[\sum_{i=0}^na_i(x)y^{(i)}=f(x),\]
in which $a_i(x)\colon\Omega_i\subseteq\bbR\to\bbR$ are called coefficients and $f(x)\colon\Omega\subseteq\bbR\to\bbR^m$ is called forcing term.
\sssc{Linear differential operator}
A linear differential operator is a polynomial of the differential operators $D$, i.e., $\dv{}{x}$ for independent variable $x$, and its order is the degree of the polynomial.
\sssc{Linearity}
Let $L$ be the linear differential operator $L[\mb{y}]=\sum_{i=0}^na_i(x)\mb{y}^{(i)}$. Then $L[\alpha\mb{y}_1+\beta\mb{y}_2]=\alpha L[\mb{y}_1]+\beta L[\mb{y}_2]$. That's why it's called to be linear.
\sssc{Homogeneity}
A linear ODE is called homogeneous iff the forcing term is 0; otherwise, it is called non-homogeneous.
\sssc{Superposition principle and solutions of homogeneous linear ODE as vector space}
If $y_1(x)$ and $y_2(x)$ are both solutions to a homogeneous linear ODE on an interval, then $c_1y_1(x)+c_2y_2(x)$ is also a solution to it on the interval for any $c_1,c_2\in\bbR$. Thus, the general solutions to a homogeneous linear ODE is a vector space.
\sssc{Complementary or associated (homogeneous) equation, natural, complementary, or homogeneous solution or function, particularly solution, and solutions of non-homogeneous linear ODE as affine space}
Given an $n$th-order linear ODE
\[\sum_{i=0}^na_i(x)y^{(i)}=f(x),\]
\[\sum_{i=0}^na_i(x)y^{(i)}=0,\]
is called its complementary or associated (homogeneous) equation, and the general solutions of the original equation is the sum of the general solutions $y_c$ of the complementary equation, called natural, complementary, or homogeneous solution or function, and a particular solution $y_p$. Thus, the general solutions to a non-homogeneous linear ODE is an affine space whose associated vector space is the vector space of the solutions to its complementary equation.
\sssc{Wrońskian or Wronskian}
For $n$ functions $f_i\colon\Omega_i\subseteq\bbR\to\bbR$ for all $i\in\bbN_{\le n}$ that are $(n-1)$ times differentiable on an interval $I$, the Wrońskian (also called Wronskian) of them, $W(f_1,f_2,\ldots,f_n)$, is a function on $I$ defined by
\[W(f_1,f_2,\ldots,f_n)(x)=\begin{vmatrix}
f_1(x) & f_2(x) & \cdots & f_n(x) \\
f_1'(x) & f_2(x) & \cdots & f_n'(x) \\
\vdots & \vdots & \ddots & \vdots \\
f_1^{(n-1)}(x) & f_2^{(n-1)}(x) & \cdots & f_n^{(n-1)}(x)
\end{vmatrix}.\]
\sssc{Non-zero Wrońskian implies linear independence}
For $n$ functions $f_i\colon\Omega_i\subseteq\bbR\to\bbR$ for all $i\in\bbN_{\le n}$ that are $(n-1)$ times differentiable on an interval $I$, if $W(f_1,f_2,\ldots,f_n)\ne0$, then they are linearly independent.
\begin{proof}
If they are not linearly independent, there exists not all zero $c_1,\ldots,c_n$ such that
\[\sum_{i=1}^nc_if_i=0\]
\[\begin{bmatrix}
f_1(x) & f_2(x) & \cdots & f_n(x) \\
f_1'(x) & f_2(x) & \cdots & f_n'(x) \\
\vdots & \vdots & \ddots & \vdots \\
f_1^{(n-1)}(x) & f_2^{(n-1)}(x) & \cdots & f_n^{(n-1)}(x)
\end{bmatrix}\begin{bmatrix}c_1\\c_2\\\vdots\\c_n\end{bmatrix}=\begin{bmatrix}0\\0\\\vdots\\0\end{bmatrix}\]
All $c_i=0$ if $W(f_1,f_2,\ldots,f_n)(x)\ne 0$ by Cramer's rule, thus $W(f_1,f_2,\ldots,f_n)(x)=0$.
\end{proof}
\sssc{General Abel's (differential equation) identity/formula}
For $n$ $(n-1)$ times differentiable solutions $y_1,y_2,\ldots,y_n$ of $y^{(n)}+\sum_{i=0}^{n-1}p_i(x)y^{(i)}=0$ on an interval $I$ where $p_i(x)$ are continuous for all $i\in[0,n-1]\cap\bbZ$, for any $x_0\in I$,
\[W(y_1,y_2,\ldots,y_n)(x)=W(y_1,y_2,\ldots,y_n)(x_0)\exp(-\int_{x_0}^xp_{n-1}(t)\dd{t}).\]
\bpr
Let $W$ denote $W(y_1,y_2,\ldots,y_n)$.
\[w(x)=\begin{vmatrix}
y_1(x) & y_2(x) & \cdots & y_n(x) \\
y_1'(x) & y_2'(x) & \cdots & y_n'(x) \\
\vdots & \vdots & \ddots & \vdots \\
y_1^{(n-1)}(x) & y_2^{(n-1)}(x) & \cdots & y_n^{(n-1)}(x)
\end{vmatrix}.\]
\[\ba
W'(x)&=\begin{vmatrix}
y_1(x) & y_2(x) & \cdots & y_n(x) \\
y_1'(x) & y_2'(x) & \cdots & y_n'(x) \\
\vdots & \vdots & \ddots & \vdots \\
y_1^{(n)}(x) & y_2^{(n)}(x) & \cdots & y_n^{(n)}(x)
\end{vmatrix}\\
&=\begin{vmatrix}
y_1(x) & y_2(x) & \cdots & y_n(x) \\
y_1'(x) & y_2'(x) & \cdots & y_n'(x) \\
\vdots & \vdots & \ddots & \vdots \\
y_1^{(n)}(x)+\sum_{i=0}^{n-2}p_i(x)y_1^{\pht{1}(i)} & y_2^{(n)}(x)+\sum_{i=0}^{n-2}p_i(x)y_2^{\pht{2}(i)} & \cdots & y_n^{(n)}(x)+\sum_{i=0}^{n-2}p_i(x)y_n^{\pht{n}(i)}
\end{vmatrix}\\
&=\begin{vmatrix}
y_1(x) & y_2(x) & \cdots & y_n(x) \\
y_1'(x) & y_2'(x) & \cdots & y_n'(x) \\
\vdots & \vdots & \ddots & \vdots \\
-p_{n-1}(x)y_1^{(n-1)}(x) & -p_{n-1}(x)y_2^{(n-1)}(x) & \cdots & -p_{n-1}(x)y_n^{(n-1)}(x)
\end{vmatrix}\\
&=-p_{n-1}(x)W(x)
\ea\]
Solve $W'(x)+p_{n-1}(x)W(x)$ using integrating factor $\exp(-\int_{x_0}^xp_{n-1}(t)\dd{t})$.
\[W(x)=W(x_0)\exp(-\int_{x_0}^xp_{n-1}(t)\dd{t}).\]
\epr
\sssc{Non-zero Wrońskian of solutions of homogeneous linear equation iff linear independence of them}
For $n$ $(n-1)$ times differentiable solutions $y_1,y_2,\ldots,y_n$ of $y^{(n)}+\sum_{i=0}^{n-1}p_i(x)y^{(i)}=0$ on an interval $I$ where $p_i(x)$ are continuous for all $i\in[0,n-1]\cap\bbZ$, for any $x_0\in I$, $W(y_1,y_2,\ldots,y_n)(x_0)\ne 0$ iff $y_1,y_2,\ldots,y_n$ are linearly independent.
\bpr
Abel's identity shows that
\[W(y_1,y_2,\ldots,y_n)(x)=W(y_1,y_2,\ldots,y_n)(x_0)\exp(-\int_{x_0}^Xp_{n-1}(t)\dd{t}).\]
$\exp(-\int_{x_0}^Xp_{n-1}(t)\dd{t})$ is always non-zero. Thus $W(y_1,y_2,\ldots,y_n)(x)$ is either zero for all $x\in I$ or non-zero for all $x\in I$.

If $y_1,y_2,\ldots,y_n$ are linearly dependent, there exists $\begin{bmatrix}c_1\\c_2\\\vdots\\c_n\end{bmatrix}\ne\begin{bmatrix}0\\0\\\vdots\\0\end{bmatrix}$ such that
\[\sum_{i=1}^nc_iy_i(x)=0\]
For any fixed $x_0\in I$, let
\[\Omega=\begin{bmatrix}
y_1(x_0) & y_2(x_0) & \cdots & y_n(x_0) \\
y_1'(x_0) & y_2'(x_0) & \cdots & y_n'(x_0) \\
\vdots & \vdots & \ddots & \vdots \\
y_1^{(n-1)}(x_0) & y_2^{(n-1)}(x_0) & \cdots & y_n^{(n-1)}(x_0)
\end{bmatrix}\]
Differentiating $n-1$ times gives
\[\Omega\begin{bmatrix}c_1\\c_2\\\vdots\\c_n\end{bmatrix}=\begin{bmatrix}0\\0\\\vdots\\0\end{bmatrix}\]
Prove by contradiction. If $\det(\Omega)=0$, by Cramer's rule, $\begin{bmatrix}c_1\\c_2\\\vdots\\c_n\end{bmatrix}=\begin{bmatrix}0\\0\\\vdots\\0\end{bmatrix}$.
\epr
\sssc{Annihilator}
An annihilator of a function $f$ is a linear differential operator $L$ such that $L(f)=0$.
\sssc{Annihilator multiplication property}
If $P$ is an annihilator of $f$, then for any linear differential operator $Q$, $PQ$ is an annihilator of $f$.
\ssc{Characteristic equation or auxiliary equation for constant coefficient homogeneous linear ODE}
\sssc{Method of complex-valued solution}
Consider an $n$th-order constant coefficient homogeneous linear ODE
\[\sum_{j=0}^na_jy^{(j)}=0,\quad a_n\neq 0,a_j\in\bbC.\]
\ben
\item The characteristic equation is
\[\sum_{j=0}^na_jr^j=0.\]
\item Suppose there are $k$ distinct roots $r_1,\ldots,r_k$, each $r_j$ with multiplicity $m_j$. (Thus, $\sum_{j=1}^km_j=n$.)
\item The $n$-parameter family of complex-valued general solution of the ODE is, with $n$ complex parameters $A_{ju}$,
\[y=\sum_{j=1}^k\sum_{u=0}^{m_j-1}A_{ju}x^ue^{r_jx}.\]
\een
\sssc{Method of real-valued solution}
Consider an $n$th-order constant coefficient homogeneous linear ODE
\[\sum_{j=0}^na_jy^{(j)}=0,\quad a_n\neq 0,a_j\in\bbR.\]
The real-valued general solution is the real part of the complex-valued solution. Alternatively, it can be solved as follows.
\ben
\item The characteristic equation is
\[\sum_{j=0}^na_jr^j=0.\]
\item Suppose there are $k$ distinct real roots $r_1,\ldots,r_k$, each $r_j$ with multiplicity $m_j$, and $h$ distinct pairs of conjugate imaginary roots $\alpha_1\pm\beta_1i,\ldots,\alpha_h\pm\beta_hi$ ($\beta_j\neq 0\forall j$), each $\alpha_j\pm\beta_ji$ with multiplicity $\omega_j$. (Thus, $\sum_{j=1}^km_j+2\sum_{j=1}^h\omega_j=n$.)
\item The $n$-parameter family of real-valued general solution of the ODE is, with $n$ real parameters $B_{ju}$, $C_{ju}$, and $D_{ju}$,
\[y=\sum_{j=1}^k\sum_{u=0}^{m_j-1}B_{ju}x^ue^{r_jx}+\sum_{j=1}^h\sum_{u=0}^{\omega_j-1}\qty(C_{ju}\cos(\beta_jx)+D_{ju}\sin(\beta_jx))x^ue^{\alpha_jx}.\]
$B_{ju}$ and $C_{ju}$ are essentially the real part of the complex parameter of complex-valued general solution for real and imaginary roots respectively, and $D_{ju}$ is essentially the imaginary part of the complex parameter of complex-valued general solution for imaginary roots.
\een
\sssc{Exponential order}
For any solution $f(t)$ of any homogeneous linear ODE of one variable $t$, there exists some $M,T\geq 0$ and $\alpha\in\bbR$ such that for all $t>T$, $\abs{f(t)}\leq e^{M\alpha t}$, and the smallest such $\alpha$ is the largest real part among all roots of its characteristic equation.
\sssc{Linear damping (阻尼)}
ODE:
\[f''(t)+2\omega_0\zeta f'(t)+\omega_0^{\pht{0}2}f(t)=\omega_0^{\pht{0}2}f(\infty)\]
where $\omega_0>0$ is called (angular) resonance/resonant frequency, $\zeta\ge0$ is called damping factor/ratio (阻尼係數/比), and $\alpha=\omega_0\zeta$ is called neper frequency (奈培頻率) or attenuation (衰減量).

Solutions:
\bit
\item $\zeta>1$ is called overdamped (過阻尼):
\[f(t)=f(\infty)+e^{-\omega_0\zeta t}\qty((f(0)-f(\infty))\cosh(\omega_0\sqrt{\zeta^2-1}t)+f'(0)\sinh(\omega_0\sqrt{\zeta^2-1}t))\]
Overshoot:
\[f(0)=f(\infty)\land f'(0)>0\implies\arg_t\max_t=\frac{\ln(\zeta+\sqrt{\zeta^2-1})}{\omega_0\sqrt{\zeta^2-1}}\]
\[f(0)=f(\infty)\land f'(0)<0\implies\arg_t\min_t=\frac{\ln(\zeta+\sqrt{\zeta^2-1})}{\omega_0\sqrt{\zeta^2-1}}\]
\item $\zeta=1$ is called critically damped (臨界阻尼):
\[f(t)=f(\infty)+\qty(f(0)-f(\infty)+f'(0)t)e^{-\omega_0t}\]
\item $0<\zeta<1$ is called underdamped (欠/輕/次阻尼):
\[f(t)=f(\infty)+e^{-\omega_0\zeta t}\qty((f(0)-f(\infty))\cos(\omega_0\sqrt{1-\zeta^2}t)+f'(0)\sin(\omega_0\sqrt{1-\zeta^2}t))\]
The damped frequency $\omega_d$ is defined as
\[\omega_d=\omega_0\sqrt{1-\zeta^2}\]
\item $\zeta=0$ is called not damped:
\[f(t)=f(\infty)+(f(0)-f(\infty))\cos(\omega_0t)+f'(0)\sin(\omega_0t)\]
\eit
Oscillation:
\bit
\item When $\zeta\ge1$, there is no oscillation:
\[(f(0)-f(\infty))f'(0)\ge0\land f(0)-f(\infty)+f'(0)\ne0\iff\forall t>0\colon (f(0)-f(\infty)+f'(0))(f(t)-k)>0\]
\item When $0\le\zeta<1$, there is oscillation:
\[(f(0)-f(\infty))f'(0)\ge0\land f(0)-f(\infty)+f'(0)\ne0\iff\exists u,v>0\colon(f(u)-f(\infty))(f(v)-f(\infty))<0\]
\eit
Assume that $\zeta>0$, then for any small enough fixed $M>0$, the smallest $t_m$ such that $t>t_m\implies\abs{f(t)-f(\infty)}<M$ is smallest at $\zeta=1$ and grows as $\zeta$ increases or decreases.
\bpr
The characteristic equation is
\[r^2+2\omega_0\zeta r+\omega_0^{\pht{0}2}=0\]
The characteristic roots $r$ are
\[r_{1,2}=-\omega_0\qty(\zeta\pm\sqrt{\zeta^2-1})\]
Case $\zeta>1$:
\[\zeta^2>\zeta^2-1>0\]
\[\zeta>\sqrt{\zeta^2-1}\]
\[\lim_{t\to\infty}e^{-\omega_0\qty(\zeta-\sqrt{\zeta^2-1})t}=0\]
\[\lim_{t\to\infty}e^{-\omega_0\qty(\zeta+\sqrt{\zeta^2-1})t}=0\]
If $f(0)=f(\infty)$:
\[f(t)=e^{-\omega_0\zeta t}f'(0)\sinh(\omega_0\sqrt{\zeta^2-1}t)\]
\[f'(t)=e^{-\omega_0\zeta t}\omega_0f'(0)\qty(\sqrt{\zeta^2-1}\cosh(\omega_0\sqrt{\zeta^2-1}t)-\zeta\sinh(\omega_0\sqrt{\zeta^2-1}t))\]
\[e^{\ln(\zeta+\sqrt{\zeta^2-1})}=\zeta+\sqrt{\zeta^2-1}\]
\[e^{-\ln(\zeta+\sqrt{\zeta^2-1})}=\frac{1}{\zeta+\sqrt{\zeta^2-1}}=\zeta-\sqrt{\zeta^2-1}\]
\[\cosh(\ln(\zeta+\sqrt{\zeta^2-1}))=\zeta\]
\[\sinh(\ln(\zeta+\sqrt{\zeta^2-1}))=\sqrt{\zeta^2-1}\]
Case $\zeta=0$:
\[r_{1,2}=\pm\omega_0i\]
\epr
\ssc{Method of undetermined coefficients (MUC) for constant coefficient non-homogeneous linear ODE}
\sssc{Superposition approach}
Consider an $n$th-order constant coefficient non-homogeneous linear ODE
\[\sum_{j=0}^na_jy^{(j)}=g(x),\quad a_n\neq 0,\]
with $g(x)$ being able to be written in the form
\[g(x)=\sum_{j=1}^pk_jx^{m_j}e^{\alpha_jx}+\sum_{j=p+1}^{p+q}x^{m_j}e^{\alpha_jx}\qty(k_j\cos(\beta_jx)+l_j\sin(\beta_jx)),\]
for some constants $p,q,m_j\in\bbN_0$ and $\alpha_j,\beta_j,k_j,l_j\in\bbR$.

The formation of its general solutions using the method of undetermined coefficients is as follows:
\ben
\item Solve the complementary equation
\[\sum_{j=0}^na_jy^{(j)}=0,\quad a_n\neq 0\]
to obtain its general solutions $y_c$.
\item For each term $k_jx^{m_j}e^{\alpha_jx}$ of $g(x)$ for $j\in\bbN\land j\leq p$, construct a particular solution $y_j$:
\ben
\item If $\alpha_j$ is not a root of the characteristic equation of the complementary equation, let $s_j=0$; otherwise, let $s_j$ be the multiplicity of $\alpha_j$ as a root of the characteristic equation of the complementary equation.
\item The ansatz is
\[y_j=x^{s_j}e^{\alpha_jx}\sum_{k=0}^{m_j}C_{jk}x^k,\]
with constants $C_{jk}\in\bbR$.
\item Solve for $C_{jk}$ such that $y_j$ is a particular solution of the constant coefficient non-homogeneous linear ODE
\[\sum_{j=0}^na_jy^{(j)}=k_jx^{m_j}e^{\alpha_jx}.\]
\een
\item For each term $x^{m_j}e^{\alpha_jx}\qty(k_j\cos(\beta_jx)+l_j\sin(\beta_jx))$ of $g(x)$ for $j\in\bbN\land p<j\leq p+q$, construct a particular solution $y_j$:
\ben
\item If $\alpha_j\pm\beta_ji$ is not a pair of conjugate imaginary roots of the characteristic equation of the complementary equation, let $s_j=0$; otherwise, let $s_j$ be the multiplicity of $\alpha_j\pm\beta_ji$ as a pair of conjugate imaginary roots of the characteristic equation of the complementary equation.
\item The ansatz is
\[y_j=x^{s_j}e^{\alpha_jx}\sum_{k=0}^{m_j}\qty(C_{jk}\cos(\beta_jx)+D_{jk}\sin(\beta_jx))x^k,\]
with constants $C_{jk},D_{jk}\in\bbR$.
\item Solve for $C_{jk},D_{jk}$ such that $y_j$ is a particular solution of the constant coefficient non-homogeneous linear ODE
\[\sum_{j=0}^na_jy^{(j)}=x^{m_j}e^{\alpha_jx}\qty(k_j\cos(\beta_jx)+l_j\sin(\beta_jx)).\]
\een
\item A particular solution of the origin ODE is
\[y_p=\sum_{j=1}^{p+q}y_j,\]
and the general solutions of the origin ODE is
\[y=y_c+y_p.\]
\een
\sssc{Annihilator approach}
Consider an $n$th-order constant coefficient non-homogeneous linear ODE
\[\sum_{j=0}^na_jy^{(j)}=g(x),\quad a_n\neq 0,\]
with $g(x)$ being able to be written in the form
\[g(x)=\sum_{j=1}^pk_jx^{m_j}e^{\alpha_jx}+\sum_{j=p+1}^{p+q}x^{m_j}e^{\alpha_jx}\qty(k_j\cos(\beta_jx)+l_j\sin(\beta_jx)),\]
for some constants $p,q,m_j\in\bbN_0$ and $\alpha_j,\beta_j,k_j,l_j\in\bbR$.

The formation of its general solutions using the method of undetermined coefficients is as follows:
\ben
\item Solve the complementary equation
\[\sum_{j=0}^na_jy^{(j)}=0,\quad a_n\neq 0\]
to obtain its general solutions $y_c$.
\item For each term $k_jx^{m_j}e^{\alpha_jx}$ of $g(x)$, a lowest-order annihilator is
\[(D-\alpha_j)^{m_j+1}\]
\bpr
Prove by Laplace transform.
\[\mathcal{L}\{k_jx^{m_j}e^{\alpha_jx}\}=\frac{k_jm_j!}{(s-\alpha_j)^{m_j+1}}\]
\epr
\item For each term $x^{m_j}e^{\alpha_jx}\qty(k_j\cos(\beta_jx)+l_j\sin(\beta_jx))$ of $g(x)$, a lowest-order annihilator is
\[(D^2-2\alpha_jD+\alpha_j^{\pht{j}2}+\beta_j^{\pht{j}2})^{m_j+1}\]
\bpr
Prove by Laplace transform.
\[\ba
\mathcal{L}\{x^{m_j}e^{\alpha_jx}\qty(k_j\cos(\beta_jx)+l_j\sin(\beta_jx))\}&=\mathcal{L}\{x^{m_j}e^{\alpha_jx}\qty((k_j+l_j)e^{i\beta_jx}-(k_j-l_j)ie^{-i\beta_jx})\}\\
&=(k_j+l_j)\mathcal{L}\{x^{m_j}e^{\alpha_jx+i\beta_jx}\}-(k_j-l_j)i\mathcal{L}\{x^{m_j}e^{\alpha_jx-i\beta_jx}\}\\
&=\frac{(k_j+l_j)m_j!}{(s-\alpha_j-i\beta_j)^{m_j+1}}-\frac{(k_j-l_j)im_j!}{(s-\alpha_j+i\beta_j)^{m_j+1}}
\ea\]
\epr
\item Find a lowest-order annihilator of $g(x)$, $L_1$, and apply it on both sides to obtain a homogeneous linear ODE
\[L_1(\sum_{j=0}^na_jy^{(j)})=0.\]
\item Solve it and obtain its general solutions.
\item Substitute the solutions into the original ODE to obtain an explicit particular solution $y_p$, and the general solutions of the origin ODE is
\[y=y_c+y_p.\]
\een
\ssc{Variation of parameters or variation of constants method of solution of non-homogeneous linear equation}
\sssc{Method of solution}
Given an $n$th-order non-homogeneous linear ODE
\[y^{(n)}+\sum_{i=0}^{n-1}p_i(x)y^{(i)}=q(x),\]
where $p_i(x)$ for all $i\in [0,n-1]\cap\bbZ$ are continuous on An interval $I$.
\ben
\item Let $y_1(x),y_2(x),\ldots,y_n(x)$ be a basis of the vector space of the solutions of the complementary equation,
\[y^{(n)}+\sum_{i=0}^{n-1}p_i(x)y^{(i)}=0.\]
\item The particular solution $y_p$ is given by:
\[y_p=\sum_{i=1}^ny_i(x)\int_{x_0}^x\frac{W_i(t)}{W(t)}\dd{t},\]
where
\bit
\item $x_0$ is an arbitrary point in $I$,
\item $W(x)$ is the Wrońskian of $y_1(x),y_2(x),\ldots,y_n(x)$,
\item $W_i(x)$ is $W(x)$ with the $i$th column replaced by $(0,0,\ldots,0,q(x))$.
\eit
\item The general solutions are given by
\[y=y_c+y_p.\]
\een
\begin{proof}
The particular solution is given by
\[y_p=\sum_{i=1}^nc_i(x)y_i(x),\]
where $c_i(x)$ are differentiable functions which are assumed to satisfy the $n-1$ conditions
\[\sum_{i=1}^nc_i'(x)y_i^{(j)}(x)=0\]
for all $j\in\bbN_0\land j\leq n-2$.

By repeated differentiation, we obtain:
\[y_p^{(j)}=\sum_{i=1}^nc_i(x)y_i^{(j)}(x)\]
for all $j\in\bbN_0\land j\leq n-1$.

One last differentiation gives
\[y_p^{(n)}=\sum_{i=1}^nc_i(x)y_i^{(n)}(x)+\sum_{i=1}^nc_i'(x)y_i^{(n-1)}(x).\]

Substituting $y_p$ into the original ODE gives:
\[y_p^{(n)}+\sum_{i=0}^{n-1}p_i(x)y_p^{(i)}=q(x),\]
\[\sum_{i=1}^nc_i(x)y_i^{(n)}(x)+\sum_{i=1}^nc_i'(x)y_i^{(n-1)}(x)+\sum_{i=0}^{n-1}p_i(x)\sum_{i=1}^nc_i(x)y_i^{(j)}(x)=q(x),\]
\[\sum_{i=1}^nc_i'(x)y_i^{(n-1)}(x)+\sum_{i=1}^nc_i(x)\qty(y_i^{(n)}(x)+\sum_{i=0}^{n-1}p_i(x)y_i^{(j)}(x))=q(x),\]
\[\sum_{i=1}^nc_i'(x)y_i^{(n-1)}(x)=q(x).\]
We obtain a system of linear equations
\[\sum_{i=1}^nc_i'(x)y_i^{(j)}(x)=0\]
for all $j\in\bbN_0\land j\leq n-2$, and
\[\sum_{i=1}^nc_i'(x)y_i^{(n-1)}(x)=q(x).\]
By Cramer's rule, we have
\[c_i'(x)=\frac{W_i(x)}{W(x)},\quad i\in\bbN\land i\leq n.\]
\end{proof}
\sssc{For second-order equation, Green's function}
For second-order equation
\[y''+p_1(x)y'+p_0(x)y=q(x).\]
Let $y_1(x),y_2(x)$ be two linearly independent solutions of its complementary equation
\[y''+p_1(x)y'+p_0(x)y=0.\]
The particular solution (or family of general solutions with the integral constants being parameters) is given by
\[y_p(x)=-y_1(x)\int_{x_0}^x\frac{y_2(t)q(t)}{y_1(t)y_2'(t)-y_2(t)y_1'(t)}\dd{t}+y_2(x)\int_{x_0}^x\frac{y_1(t)q(t)}{y_1(t)y_2'(t)-y_2(t)y_1'(t)}\dd{t}.\]
This can be written as
\[y_p(x)=\int_{x_0}^xG(x,t)q(t)\dd{t},\]
where
\[G(x,t)=\frac{y_1(t)y_2(x)-y_1(x)y_2(t)}{y_1(t)y_2'(t)-y_2(t)y_1'(t)}\]
and is called the Green's function for the ODE or the Green's function for the linear differential operator
\[D^2+p_1(x)D+p_0(x).\]
\ssc{Reduction of order}
\sssc{Method}
Given an $n$th-order homogeneous linear ODE
\[\sum_{i=0}^np_i(x)y^{(i)}=0,\]
if a nontrivial particular solution $y_1(x)$ has been known, we can reduce the order of the ODE by 1 by the substitutions
\[y=vy_1\]
\[u=v'\]
to an $(n-1)$th-order linear ODE.
\sssc{Closed-form formula for second-order homogeneous linear ODEs}
Given a second-order linear ODE
\[y''+p(x)y'+q(x)y=0\]
and a nontrivial particular solution $y_1(x)$, then the substitutions
\[y=vy_1\]
\[u=v'\]
gives
\[y_1u'+(2y_1'+p(x)y_1)u=0\]
and another particular solution of the ODE is
\[y_2=y_1\int\frac{e^{-\int p(x)\dd{x}}}{y_1(x)^2}\dd{x}\]
and the general solution is
\[y=c_1y_1+c_2y_2\]
\begin{proof}
\[y=vy_1\]
\[y'=v'y_1+vy_1'\]
\[y''=v''y_1+2v'y_1'+vy_1''\]
\[y_1''=-py_1'-qy_1\]
\[y''=(v''-qv)y_1+(2v'-pv)y_1'\]
\[y''=-py'-qy\]
\[y''=(-pv'-qv)y_1-pvy_1'\]
\[(v''+pv')y_1+2v'y_1'=0\]
\[y_1v''+(2y_1'+py_1)v'=0\]
\[u=v'\]
\[y_1u'+(2y_1'+py_1)u=0\]
\[u'+\qty(\frac{2y_1'}{y_1}+p)u=0\]
\[\mu(x)=\exp(\int\frac{2y_1'}{y_1}+p\dd{x})=y_1^{\phantom{1}2}e^{\int p\dd{x}}\]
\[(u\mu)'=0\]
\[u=\frac{c_2}{\mu}=\frac{c_2e^{-\int p\dd{x}}}{y_1^{\phantom{1}2}}\]
\[v=\int u\dd{x}=c_2\int\frac{e^{-\int p\dd{x}}}{y_1^{\phantom{1}2}}\dd{x}+c_1\]
\[y=vy_1=c_2y_1\int\frac{e^{-\int p\dd{x}}}{y_1^{\phantom{1}2}}\dd{x}+c_1y_1\]
\[y_2=y_1\int\frac{e^{-\int p(x)\dd{x}}}{y_1(x)^2}\dd{x}\]
\[y=c_1y_1+c_2y_2\]
\end{proof}
\sssc{Closed-form formula for second-order homogeneous linear ODEs}
Given a second-order linear ODE
\[y''+p(x)y'+q(x)y=g(x)\]
and a nontrivial solution of the associated homogeneous equation $y_1(x)$, then the substitutions
\[y=vy_1\]
\[u=v'\]
gives
\[y_1u'+(2y_1'+p(x)y_1)u=g(x)\]
and another solution of the associated homogeneous equation is
\[y_2=y_1\int\frac{e^{-\int p(x)\dd{x}}}{y_1(x)^2}\dd{x}\]
and a particular solution is
\[y_p=y_1\int\frac{e^{-\int p\dd{x}}}{y_1^{\phantom{1}2}}\int g(x)y_1e^{\int p\dd{x}}\dd{x}\dd{x}\]
and the general solution is
\[y=c_1y_1+c_2y_2+y_p\]
\begin{proof}
\[y=vy_1\]
\[y'=v'y_1+vy_1'\]
\[y''=v''y_1+2v'y_1'+vy_1''\]
\[y_1''=-py_1'-qy_1\]
\[y''=(v''-qv)y_1+(2v'-pv)y_1'\]
\[y''=-py'-qy+g\]
\[y''=(-pv'-qv)y_1-pvy_1'+g\]
\[(v''+pv')y_1+2v'y_1'=g\]
\[y_1v''+(2y_1'+py_1)v'=g\]
\[u=v'\]
\[y_1u'+(2y_1'+py_1)u=g\]
\[u'+\qty(\frac{2y_1'}{y_1}+p)u=\frac{g}{y_1}\]
\[\mu(x)=\exp(\int\frac{2y_1'}{y_1}+p\dd{x})=y_1^{\phantom{1}2}e^{\int p\dd{x}}\]
\[(u\mu)'=\frac{g\mu}{y_1}\]
\[u=\frac{\int\frac{g\mu}{y_1}\dd{x}}{\mu}=\frac{e^{-\int p\dd{x}}}{y_1^{\phantom{1}2}}\int g(x)y_1e^{\int p\dd{x}}\dd{x}\]
\[v=\int u\dd{x}=\int\frac{e^{-\int p\dd{x}}}{y_1^{\phantom{1}2}}\int g(x)y_1e^{\int p\dd{x}}\dd{x}\dd{x}+c_2\int\frac{e^{-\int p\dd{x}}}{y_1^{\phantom{1}2}}\dd{x}+c_1\]
\[y=vy_1=y_1\int\frac{e^{-\int p\dd{x}}}{y_1^{\phantom{1}2}}\int g(x)y_1e^{\int p\dd{x}}\dd{x}\dd{x}+c_2y_1\int\frac{e^{-\int p\dd{x}}}{y_1^{\phantom{1}2}}\dd{x}+c_1y_1\]
\[y_2=y_1\int\frac{e^{-\int p(x)\dd{x}}}{y_1(x)^2}\dd{x}\]
\[y_p=y_1\int\frac{e^{-\int p\dd{x}}}{y_1^{\phantom{1}2}}\int g(x)y_1e^{\int p\dd{x}}\dd{x}\dd{x}\]
\[y=c_1y_1+c_2y_2+y_p\]
\end{proof}
\ssc{(Cauchy-)Euler or equidimensional equation}
\sssc{Definition} 
An $n$th order Cauchy-Euler equation has the form
\[\sum_{i=0}^na_ix^iy^{(i)}=g(x),\quad a_n\ne 0.\]
\sssc{Method of solution}
Consider an $n$th-order Cauchy-Euler equation
\[\sum_{i=0}^na_ix^iy^{(i)}=g(x),\quad a_n\ne 0.\]
The formation of its general solutions using the characteristic equation (also called auxiliary equation) is as follows:
\ben
\item The complementary equation is
\[\sum_{i=0}^na_ix^iy^{(i)}=0\]
\item The characteristic equation is
\[\sum_{j=0}^na_jr^j=0.\]
\item Suppose there are $k$ distinct real roots $r_1,r_2,\ldots,r_k$, each $r_j$ with multiplicity $m_j$, and $h$ distinct pairs of conjugate imaginary roots $\alpha_1\pm\beta_1i,\alpha_2\pm\beta_2i,\ldots,\alpha_h\pm\beta_hi$ ($\beta_j\neq 0$), each $\alpha_j\pm\beta_ji$ with multiplicity $\omega_j$. (Thus, $\sum_{j=1}^km_j+2\sum_{j=1}^h\omega_j=n$.)
\item The $n$-parameter family of solutions of the complementary equation is
\[y_c=\sum_{j=1}^k\sum_{u=0}^{m_j-1}A_{ju}(\ln x)^ux^{r_j}+\sum_{j=1}^h\sum_{u=0}^{\omega_j-1}\qty(B_{ju}\cos(\beta_j\ln x)+C_{ju}\sin(\beta_j\ln x))(\ln x)^ux^{\alpha_j},\]
with parameters $A_{ju},B_{ju},C_{ju}$.
\item Then, we use variation of parameters. Let $y_1(x),y_2(x),\ldots,y_n(x)$ be a basis of the vector space of $y_c$.

The particular solution $y_p$ is given by:
\[y_p=\sum_{i=1}^ny_i(x)\int_{x_0}^x\frac{W_i(t)}{W(t)}\dd{t},\]
where
\bit
\item $x_0$ is an arbitraty point in the interval of interest,
\item $W(x)$ is the Wrońskian of $y_1(x),y_2(x),\ldots,y_n(x)$,
\item $W_i(x)$ is the Wrońskian of $W(x)$ but with the $i$th column replaced by $(0,0,\ldots,0,g(x))$.
\eit
\item The general solutions are given by
\[y=y_c+y_p.\]
\een
\ssc{Power series solutions and Frobenius method for second-order homogeneous linear ODE}
\sssc{Type of Points}
Consider a second-order homogeneous linear ODE
\[f(x)y''+g(x)y'+h(x)y=0\]
where $f,g,h$ are continuous on interval $I$.

Let
\[p(x)=\frac{g(x)}{f(x)},\quad q(x)=\frac{h(x)}{f(x)}.\]
For any $a\in I$, it is called an ordinary point iff $p(x)$ and $q(x)$ are both analytic at $a$ and called a singular point otherwise.

For a singular point $a$, it is called a regular singular point iff $(x-a)p(x)$ and $(x-a)^2q(x)$ are both analytic at $a$ and called an irregular singular point otherwise.
\sssc{Power series solution or solution about ordinary points}
Consider a second-order homogeneous linear ODE
\[f(x)y''+g(x)y'+h(x)y=0\]
where $f,g,h$ are continuous on interval $I$.

Let
\[p(x)=\frac{g(x)}{f(x)},\quad q(x)=\frac{h(x)}{f(x)}.\]
\ben
\item We seek a power series solution of
\[y''+p(x)y'+q(x)y=0\]
in the form:
\[y=\sum_{n=0}^{\infty}a_n(x-x_0)^n,\]
where $x_0$ is a fixed ordinary point called expansion point and $a_n$ are constants to be determined. There are always two linearly independent solutions of this form.
\item Differentiate the series term by term:
\[y'=\sum_{n=0}^{\infty}a_nn(x-x_0)^{n-1}=\sum_{n=0}^{\infty}a_{n+1}(n+1)(x-x_0)^n,\]
\[y''=\sum_{n=0}^{\infty}a_nn(n-1)(x-x_0)^{n-2}=\sum_{n=0}^{\infty}a_{n+2}(n+2)(n+1)(x-x_0)^n.\]
\item Substitute into the ODE, rewrite $p(x)$ and $q(x)$ into Taylor expansion at $x_0$, rewrite the infinite series products (if any) into Cauchy products, and rewrite the LHS as a power series in $x^n$.
\item Solve for the recurrence relation of $a_n$. We may set $a_0$ as any nonzero constant, typically 1.
\item Let $R$ be the distance between $x_0$ and the nearest singular point if there is one or $\infty$ if there is none. The two-parameter family of general solutions for all $0<|x-x_0|<R$ with parameters $a_0,a_1$ is:
\[y=\sum_{n=0}^{\infty}a_n(x-x_0)^n.\]
\een
\bpr
TODO
\epr
\sssc{Frobenius method or solution about singular points}
Consider a second-order homogeneous linear ODE
\[f(x)y''+g(x)y'+h(x)y=0.\]
Let
\[p(x)=\frac{g(x)}{f(x)},\quad q(x)=\frac{h(x)}{f(x)}.\]
\ben
\item We seek a solution, called Frobenius series, of
\[(x-x_0)^2y''+(x-x_0)r(x)y'+s(x)y=0\]
in the form:
\[y=\sum_{n=0}^{\infty}a_n(x-x_0)^{n+r},\quad a_0\neq 0,\]
where $x_0$ is a fixed regular singular point called expansion point, $r(x)=(x-x_0)p(x)$ and $s(x)=(x-x_0)^2q(x)$ are analytic at $x_0$, and $a_n$ are constants to be determined.
\item Differentiate the series term by term:
\[y'=\sum_{n=0}^{\infty}a_n(n+r)(x-x_0)^{n+r-1},\]
\[y''=\sum_{n=0}^{\infty}a_n(n+r)(n+r-1)(x-x_0)^{n+r-2}.\]
\item Substitute into the ODE, rewrite $(x-x_0)r(x)$ and $s(x)$ into Taylor expansion at $x_0$, rewrite the infinite series products (if any) into Cauchy products, and let the sum of all terms involving $a_0$ of the three be $L$.
\item The indicial polynomial $I(r)$ is defined as the coefficient of the lowest power term of $\frac{L}{a_0}$, and the indicial equation is defined as
\[I(r)=0,\]
which is a quadratic equation for $r$.
\item Solve the indicial equation for $r$. Let the roots be $r_1\geq r_2$, called indical roots or exponents.
\item Let $R$ be the distance between $x_0$ and the nearest singular point if there is one or $\infty$ if there is none. Depending on the relation of $r_1$ and $r_2$:
\bit
\item Case $r_1-r_2\notin\bbZ\land r_1,r_2\in\bbR$: Two linearly independent solutions for all $0<|x-x_0|<R$ are:
\[y_1(x)=(x-x_0)^{r_1}\sum_{n=0}^{\infty}A_n(x-x_0)^n,\quad A_0\neq 0\land A_n\in\bbR,\]
\[y_2(x)=(x-x_0)^{r_2}\sum_{n=0}^{\infty}B_n(x-x_0)^n,\quad B_0\neq 0\land B_n\in\bbR.\]
\item Case $r_1,r_2=\alpha\pm\beta i,\quad\beta\ne0$: Two linearly independent solutions for all $0<|x-x_0|<R$ are:
\[y_1(x)=(x-x_0)^{\alpha}\cos(\beta\ln|x-x_0|)\sum_{n=0}^{\infty}A_n(x-x_0)^n,\quad A_0\neq 0\land A_n\in\bbR,\]
\[y_2(x)=(x-x_0)^{\alpha}\sin(\beta\ln|x-x_0|)\sum_{n=0}^{\infty}B_n(x-x_0)^n,\quad B_0\neq 0\land B_n\in\bbR.\]
\item Case $r_1-r_2\in\bbN_0$: Two solutions for all $0<|x-x_0|<R$ are:
\[y_1(x)=(x-x_0)^{r_1}\sum_{n=0}^{\infty}A_n(x-x_0)^n,\quad A_0\neq 0\land A_n\in\bbR,\]
\[y_2(x)=(x-x_0)^{r_2}\sum_{n=0}^{\infty}B_n(x-x_0)^n,\quad B_0\neq 0\land B_n\in\bbR.\]
If they are linearly independent, no modification is needed; otherwise, $y_2$ is modified to
\[y_2(x)=\ln|x-x_0|y_1(x)+(x-x_0)^{r_2}\sum_{n=0}^{\infty}B_n(x-x_0)^n,\quad B_0\neq 0\land B_n\in\bbR,\]
where $B_n$ is different from the original one.
\eit 
\item Substitute into the ODE, rewrite $(x-x_0)r(x)$ and $s(x)$ into Taylor expansion at $x_0$, rewrite the infinite series products (if any) into Cauchy products, and rewrite the LHS as a power series in $x^n$.
\item Solve for the recurrence relations of $A_n$ and $B_n$.
\bit
\item Case $(r_1-r_2\notin\bbZ\land r_1,r_2\in\bbR)\lor (r_1,r_2=\alpha\pm\beta i,\quad\beta\ne0)$: We may set $A_0$ and $B_0$ as any nonzero constants (typically 1).
\item Case $r_1-r_2\in\bbN_0$: The terms multiplied by $\ln|x-x_0|$ in the substitution of $y_2(x)$ yield the same recurrence relation for $A_n$ as that obtained from the substitution of $y_1(x)$, and therefore need not be solved again. We may set $A_0$ as any nonzero constant (typically 1). The lowest two-order equations in the substitution of $y_2(x)$ may determine $B_0$ and/or $B_1$, or relation of $B_0$ and $B_1$. If one of them remains free after applying the constraint, it may be chosen arbitrarily with the constraint that $B_0$ can't be 0, and the remaining can be determined uniquely by the recurrence relation.
\eit
\item The two-parameter family of general solutions for all $0<|x-x_0|<R$ is:
\[y=C_1y_1(x)+C_2y_2(x),\quad C_1,C_2\in\bbR.\]
\een
\begin{proof}
TODO
\end{proof}
\ssc{Laplace transform}
\sssc{Constant coefficients}
To solve
\[\sum_{i=0}^na_iy^{(i)}=g(x)\]
assuming that $y$ is of exponential order.

Laplace transform both sides:
\[\sum_{j=0}^na_j\qty(s^jY(s)-\sum_{i=0}^{j-1}s^{j-1-i}\dv[i]{y}{t}(0))=G(s)\]
Let
\[P(s)=\sum_{j=0}^na_js^j\]
\[Q(s)=\sum_{j=0}^na_j\sum_{i=0}^{j-1}s^{j-1-i}\dv[i]{y}{t}(0)\]
Thus
\[P(s)Y(s)=Q(s)+G(s),\]
where \(P(s)\) and \(Q(s)\) are polynomials of $s$ of degrees $n$ and $<n$.

We call
\[W(s)=\frac{1}{P(s)}\]
transfer function of the system and write
\[Y(s)=W(s)Q(s)+W(s)G(s)\]
\[y_0(t)=\mathcal{L}^{-1}\{W(s)Q(s)\}\]
\[y_1(t)=\mathcal{L}^{-1}\{W(s)G(s)\}\]
\[y(t)=y_0(t)+y_1(t)\]
We call $y_0(t)$ natural, homogeneous, unforced, or zero-input response (ZIR) and $y_1(t)$ forced, driven, particular, or zero-state responses (ZSR).
\ssc{Bessel functions and Legendre polynomials}
\sssc{Bessel functions}
Bessel functions are a family of special functions that are solutions to the Bessel's equation:
\[x^2y''+xy'+(x^2-\nu^2)y=0,\]
where $x\in\bbR$ is the independent variable and $\nu\in\bbR$ is called order.
\bit
\item Bessel function of the first kind:
    \[J_{\nu}(x)=\sum_{m=0}^{\infty}\lim_{\mu\to\nu}\frac{(-1)^m}{m!\Gamma(m+\mu+1)}\qty(\frac{x}{2})^{2m+\nu}.\]
\bit
\item For $\nu\in\bbR\setminus\bbZ$:
    \[J_{\nu}(x)=\sum_{m=0}^{\infty}\frac{(-1)^m}{m!\Gamma(m+\nu+1)}\qty(\frac{x}{2})^{2m+\nu}.\]
    $J_{\nu}(x)$ and $J_{-\nu}(x)$ are two linearly independent solutions of the Bessel's equation. $J_{\nu}(x)$ is defined on $\bbR_{>=0}$ for $\nu\in\bbR_{>0}\setminus\bbZ$, and on $\bbR_{>0}$ for $\nu\in\bbR_{<0}\setminus\bbZ$. $J_{\nu}(0)=0$ for $\nu\in\bbR_{>0}\setminus\bbZ$.
\item For $\nu\in\bbZ$:
    \[J_{\nu}(x)=\sum_{m=0}^{\infty}\lim_{\mu\to\nu}\frac{(-1)^m}{m!\Gamma(m+\mu+1)}\qty(\frac{x}{2})^{2m+\nu},\]
    that is,
    \[J_{\abs{\nu}}(x)=\sum_{m=0}^{\infty}\frac{(-1)^m}{m!(m+\abs{\nu})!}\qty(\frac{x}{2})^{2m+\abs{\nu}}\]
    \[J_{-\abs{\nu}}(x)=\sum_{m=0}^{\infty}\frac{(-1)^{m+\abs{\nu}}}{m!(m+\abs{\nu})!}\qty(\frac{x}{2})^{2m+\abs{\nu}},\]
    that is,
\bit
\item For $\nu\in\bbN_0$:
\[J_{\nu}(x)=\sum_{m=0}^{\infty}\frac{(-1)^m}{m!(m+\nu)!}\qty(\frac{x}{2})^{2m+\nu}\]
\[J_{-\nu}(x)=\sum_{m=0}^{\infty}\frac{(-1)^{m+\nu}}{m!(m+\nu)!}\qty(\frac{x}{2})^{2m+\nu}\]
\item For $-\nu\in\bbN_0$:
\[J_{\nu}(x)=\sum_{m=0}^{\infty}\frac{(-1)^{m-\nu}}{m!(m-\nu)!}\qty(\frac{x}{2})^{2m-\nu}\]
\[J_{-\nu}(x)=\sum_{m=0}^{\infty}\frac{(-1)^m}{m!(m-\nu)!}\qty(\frac{x}{2})^{2m-\nu}\]
\eit
and
\[J_{-\nu}(x)=J_{\nu}(-x)=(-1)^{\nu}J_{\nu}(x)\]
$J_{\nu}(x)$ is defined on $\bbR$. $J_0(0)=1$. $J_{\nu}(0)=0$ for $\nu\in\bbZ_{\ne0}$.
\eit
\item Bessel function of the second kind:
    \[Y_{\nu}(x)=\lim_{\mu\to\nu}\frac{J_{\mu}(x)\cos(\mu\pi)-J_{-\mu}(x)}{\sin(\mu\pi)}.\]
\bit
\item For $\nu\in\bbR\setminus\bbZ$:
    \[Y_{\nu}(x)=\frac{J_{\nu}(x)\cos(\nu\pi)-J_{-\nu}(x)}{\sin(\nu\pi)},\]
    which is a linear combination of $J_{\nu}(x)$ and $J_{-\nu}(x)$.
\item For $\nu\in\bbZ$:
    \[Y_{\nu}(x)=\lim_{\mu\to\nu}Y_{\mu}(x),\]
    which is a solution that is linearly independent to $J_{\nu}(x)$. $Y_{\nu}(x)$ is defined on $\bbR_{>0}$.
    \[\lim_{x\to0^+}Y_{\nu}(x)=-\infty.\]
    \[Y_0(x)=\frac{2}{\pi}J_0(x)\qty(\gamma+\ln(\frac{x}{2}))-\frac{2}{\pi}\sum_{k=1}^{\infty}\frac{(-1)^k}{(k!)^2}\qty(sum_{j=1}^k\frac{1}{j})\qty(\frac{x}{2})^{2k}.\]
\eit
\eit
The general solution is:
\bit
\item For $\nu\in\bbR$:
    \[y(x)=CJ_{\nu}(x)+DY_{\nu}(x).\]
\item For $\nu\in\bbR\setminus\bbZ$:
    \[y(x)=CJ_{\nu}(x)+DJ_{-\nu}(x).\]
\eit
\bpr
\[x^2y''+xy'+(x^2-\nu^2)y=0\]
\[y''+x^{-1}y'+(1-\nu^2x^{-2})y=0\]
$0$ is a regular singular point.
\[y=\sum_{n=0}^{\infty}a_nx^{n+r}\]
\[y'=\sum_{n=0}^{\infty}a_n(n+r)x^{n+r-1}\]
\[y''=\sum_{n=0}^{\infty}a_n(n+r)(n+r-1)x^{n+r-2}\]
\[\sum_{n=0}^{\infty}a_n(n+r)(n+r-1)x^{n+r}+\sum_{n=0}^{\infty}a_n(n+r)x^{n+r}+\sum_{n=0}^{\infty}a_nx^{n+r+2}-\nu^2\sum_{n=0}^{\infty}a_nx^{n+r}=0\]
\[L=a_0r(r-1)x^r+a_0rx^r+a_0x^{r+2}-\nu^2a_0x^r\]
\[I(r)=r^2-\nu^2=0\]
\[r=\pm\nu\]
\[y_1(x)=\sum_{n=0}^{\infty}A_nx^{n+\abs{\nu}}\]
\[y_1'(x)=\sum_{n=0}^{\infty}A_n(n+\abs{\nu})x^{n+\abs{\nu}-1}\]
\[y_1''(x)=\sum_{n=0}^{\infty}A_n(n+\abs{\nu})(n+\abs{\nu}-1)x^{n+\abs{\nu}-2}\]
\[y_2(x)=\sum_{n=0}^{\infty}B_nx^{n-\abs{\nu}}\]
\[y_2'(x)=\sum_{n=0}^{\infty}B_n(n-\abs{\nu})x^{n-\abs{\nu}-1}\]
\[y_2''(x)=\sum_{n=0}^{\infty}B_n(n-\abs{\nu})(n-\abs{\nu}-1)x^{n-\abs{\nu}-2}\]
\[\sum_{n=0}^{\infty}A_n(n+\abs{\nu})(n+\abs{\nu}-1)x^{n+\abs{\nu}}+\sum_{n=0}^{\infty}A_n(n+\abs{\nu})x^{n+\abs{\nu}}+\sum_{n=0}^{\infty}A_nx^{n+\abs{\nu}+2}-\abs{\nu}^2\sum_{n=0}^{\infty}A_nx^{n+\abs{\nu}}=0\]
\[A_1(1+2\abs{\nu})x^{1+\abs{\nu}}+\sum_{n=2}^{\infty}\qty(A_n(n^2+2n\abs{\nu})+A_{n-2})x^{n+\abs{\nu}}=0\]
\[A_n=0,\quad\frac{n+1}{2}\in\bbN\]
\[A_n(n^2+2n\abs{\nu})+A_{n-2}=0,\quad\frac{n}{2}\in\bbN\]
\[A_{2m}=\frac{(-1)^mA_0}{\prod_{i=1}^m(2i)(2i+2\abs{\nu})}=\lim_{\mu\to\abs{\nu}}\frac{(-1)^mA_0\Gamma(\mu+1)}{2^{2m}m!\Gamma(m+\mu+1)},\quad m\in\bbN\]
Let
\[A_0=\lim_{\mu\to\abs{\nu}}\frac{1}{2^\mu\Gamma(\mu+1)}\]
\[A_{2m}=\lim_{\mu\to\abs{\nu}}\frac{(-1)^m}{2^{2m+\mu}m!\Gamma(m+\mu+1)},\quad m\in\bbN\]
\[y_1(x)=J_{\abs{\nu}}(x)=\sum_{m=0}^{\infty}\lim_{\mu\to\abs{\nu}}\frac{(-1)^m}{2^{2m+\mu}m!\Gamma(m+\mu+1)}x^{2m+\abs{\nu}}=\sum_{m=0}^{\infty}\lim_{\mu\to\abs{\nu}}\frac{(-1)^m}{m!\Gamma(m+\mu+1)}\qty(\frac{x}{2})^{2m+\abs{\nu}}\]
\[y_2(x)=J_{-\abs{\nu}}(x)=\sum_{m=0}^{\infty}\lim_{\mu\to\abs{\nu}}\frac{(-1)^m}{m!\Gamma(m-\mu+1)}\qty(\frac{x}{2})^{2m-\abs{\nu}}\]
For $\nu\in\bbR\setminus\bbZ$, that is,
\[J_{\nu}(x)=\sum_{m=0}^{\infty}\frac{(-1)^m}{m!\Gamma(m+\nu+1)}\qty(\frac{x}{2})^{2m+\nu}\]
\[J_{-\nu}(x)=\sum_{m=0}^{\infty}\frac{(-1)^m}{m!\Gamma(m-\nu+1)}\qty(\frac{x}{2})^{2m-\nu}\]
For $2\nu\in\bbZ$, test whether they are are linearly dependent.
\[J_{\nu}'(x)=\sum_{m=0}^{\infty}\lim_{\mu\to\nu}\frac{(-1)^m}{m!\Gamma(m+\mu+1)}\qty(\frac{x}{2})^{2m+\nu-1}\frac{2m+\nu}{2}\]
\[J_{-\nu}'(x)=\sum_{m=0}^{\infty}\lim_{\mu\to\nu}\frac{(-1)^m}{m!\Gamma(m-\mu+1)}\qty(\frac{x}{2})^{2m-\nu-1}\frac{2m-\nu}{2}\]
\[W(J_{\nu},J_{-\nu})(x)=\sum_{m=0}^{\infty}\sum_{n=0}^{\infty}\lim_{\mu\to\nu}\frac{(-1)^{m+n}(m-n-\nu)}{m!n!\Gamma(m-\mu+1)\Gamma(n+\mu+1)}\qty(\frac{x}{2})^{2m+2n-1}\]
The coefficient of the $x^{-1}$ term is
\[\lim_{\mu\to\nu}\frac{-2\nu}{\Gamma(1-\mu)\Gamma(1+\mu)}=\lim_{\mu\to\nu}\frac{-2}{\Gamma(\mu)\Gamma(1-\mu)}=\frac{-2\sin(\pi\nu)}{\pi}\]
Thus for $\nu\in\bbZ$, $W(J_{\nu},J_{-\nu})=0$,
\[J_{\abs{\nu}}(x)=\sum_{m=0}^{\infty}\frac{(-1)^m}{m!(m+\abs{\nu})!}\qty(\frac{x}{2})^{2m+\abs{\nu}}\]
and observe that for $k\in\bbN_0$,
\[\lim_{\mu\to\nu}\frac{1}{\Gamma(-k)}=0,\]
thus
\[J_{-\abs{\nu}}(x)=\sum_{m=0}^{\infty}\frac{(-1)^{m+\abs{\nu}}}{m!(m+\abs{\nu})!}\qty(\frac{x}{2})^{2m+\abs{\nu}},\]
and
\[J_{-\abs{\nu}}(x)=(-1)^{\abs{\nu}}J_{\abs{\nu}}(x),\]
that is,
\bit
\item For $\nu\in\bbN_0$:
\[J_{\nu}(x)=\sum_{m=0}^{\infty}\frac{(-1)^m}{m!(m+\nu)!}\qty(\frac{x}{2})^{2m+\nu}\]
\[J_{-\nu}(x)=\sum_{m=0}^{\infty}\frac{(-1)^{m+\nu}}{m!(m+\nu)!}\qty(\frac{x}{2})^{2m+\nu}\]
\[J_{-\nu}(x)=(-1)^{\nu}J_{\nu}(x)\]
\item For $-\nu\in\bbN_0$:
\[J_{\nu}(x)=\sum_{m=0}^{\infty}\frac{(-1)^{m-\nu}}{m!(m-\nu)!}\qty(\frac{x}{2})^{2m-\nu}\]
\[J_{-\nu}(x)=\sum_{m=0}^{\infty}\frac{(-1)^m}{m!(m-\nu)!}\qty(\frac{x}{2})^{2m-\nu}\]
\[J_{-\nu}(x)=(-1)^{\nu}J_{\nu}(x)\]
\eit
and a new $y_2(x)$ is
\[y_2(x)=\ln|x|J_{\abs{\nu}}(x)+\sum_{n=0}^{\infty}B_nx^{n-\abs{\nu}}\]
\[y_2'(x)=\ln|x|J_{\abs{\nu}}'(x)+x^{-1}J_{\abs{\nu}}(x)+\sum_{n=0}^{\infty}B_n(n-\abs{\nu})x^{n-\abs{\nu}-1}\]
\[y_2''(x)=\ln|x|J_{\abs{\nu}}''(x)+2x^{-1}J_{\abs{\nu}}'(x)-x^{-2}J_{\abs{\nu}}(x)+\sum_{n=0}^{\infty}B_n(n-\abs{\nu})(n-\abs{\nu}-1)x^{n-\abs{\nu}-2}\]
\[x^2\ln|x|J_{\abs{\nu}}''(x)+2xJ_{\abs{\nu}}'(x)-J_{\abs{\nu}}(x)+\sum_{n=0}^{\infty}B_n(n-\abs{\nu})(n-\abs{\nu}-1)x^{n-\abs{\nu}}+x\ln|x|J_{\abs{\nu}}'(x)+J_{\abs{\nu}}(x)+\sum_{n=0}^{\infty}B_n(n-\abs{\nu})x^{n-\abs{\nu}}+x^2\ln|x|J_{\abs{\nu}}(x)+\sum_{n=0}^{\infty}B_nx^{n-\abs{\nu}+2}-\abs{\nu}^2\ln|x|J_{\abs{\nu}}(x)-\abs{\nu}^2\sum_{n=0}^{\infty}B_nx^{n-\abs{\nu}}=0\]
\[2xJ_{\abs{\nu}}'(x)+B_1(1-2\abs{\nu})x^{1-\abs{\nu}}+\sum_{n=2}^{\infty}(B_nn(n-2\abs{\nu})+B_{n-2})x^{n-\abs{\nu}}=0\]
\[J_{\abs{\nu}}'(x)=\sum_{m=0}^{\infty}\frac{(-1)^m(2m+\abs{\nu})}{2m!(m+\abs{\nu})!}\qty(\frac{x}{2})^{2m+\abs{\nu}-1}\]
\[2xJ_{\abs{\nu}}'(x)=\sum_{m=0}^{\infty}\frac{(-1)^m2(2m+\abs{\nu})}{m!(m+\abs{\nu})!}\qty(\frac{x}{2})^{2m+\abs{\nu}}\]
\[\sum_{m=0}^{\infty}\frac{(-1)^m2(2m+\abs{\nu})}{m!(m+\abs{\nu})!}\qty(\frac{x}{2})^{2m+\abs{\nu}}+B_1(1-2\abs{\nu})x^{1-\abs{\nu}}+\sum_{n=2}^{\infty}(B_nn(n-2\abs{\nu})+B_{n-2})x^{n-\abs{\nu}}=0\]
TODO
How this = Y
\epr
\[J_{\frac{1}{2}}(x)=\frac{2}{\pi x}\sin(x)\]
\[J_{-\frac{1}{2}}(x)=\frac{2}{\pi x}\cos(x)\]
\[Y_{\frac{1}{2}}(x)=-\frac{2}{\pi x}\cos(x)\]
\[Y_{-\frac{1}{2}}(x)=\frac{2}{\pi x}\sin(x)\]
\bpr
\[\Gamma(\frac{1}{2})=\sqrt{\pi}\]
\[\Gamma(m+\frac{1}{2}+1)=\prod_{i=0}^m\frac{2i+1}{2}\sqrt{\pi}=\frac{(2m+1)!}{2^{2m+1}m!}\sqrt{\pi}\]
\[\ba
J_{\frac{1}{2}}(x)&=\sum_{m=0}^{\infty}\frac{(-1)^m}{m!\Gamma(m+\frac{1}{2}+1)}\qty(\frac{x}{2})^{2m+\frac{1}{2}}\\
&=\sum_{m=0}^{\infty}\frac{(-1)^m2^{2m+1}\sqrt{\pi}}{(2m+1)!}\qty(\frac{x}{2})^{2m+\frac{1}{2}}\\
&=\sqrt{\frac{2}{\pi x}}\sum_{m=0}^{\infty}\frac{(-1)^m}{(2m+1)!}x^{2m+1}\\
&=\sqrt{\frac{2}{\pi x}}\sin(x)
\ea\]
\[\Gamma(m-\frac{1}{2}+1)=\prod_{i=1}^m\frac{2i-1}{2}\sqrt{\pi}=\frac{(2m)!}{2^{2m}m!}\sqrt{\pi}\]
\[\ba
J_{-\frac{1}{2}}(x)&=\sum_{m=0}^{\infty}\frac{(-1)^m}{m!\Gamma(m-\frac{1}{2}+1)}\qty(\frac{x}{2})^{2m-\frac{1}{2}}\\
&=\sum_{m=0}^{\infty}\frac{(-1)^m2^{2m}\sqrt{\pi}}{(2m)!}\qty(\frac{x}{2})^{2m-\frac{1}{2}}\\
&=\sqrt{\frac{2}{\pi x}}\sum_{m=0}^{\infty}\frac{(-1)^m}{(2m)!}x^{2m}\\
&=\sqrt{\frac{2}{\pi x}}\cos(x)
\ea\]
\[Y_{\frac{1}{2}}(x)=\frac{J_{\frac{1}{2}}(x)\cos(\frac{1}{2}\pi)-J_{-\frac{1}{2}}(x)}{\sin(\frac{1}{2}\pi)}=-J_{-\frac{1}{2}}(x)=-\frac{2}{\pi x}\cos(x)\]
\[Y_{-\frac{1}{2}}(x)=\frac{J_{-\frac{1}{2}}(x)\cos(-\frac{1}{2}\pi)-J_{\frac{1}{2}}(x)}{\sin(-\frac{1}{2}\pi)}=J_{\frac{1}{2}}(x)=\frac{2}{\pi x}\sin(x)\]
\epr
\sssc{Modified Bessel functions}
Modified Bessel functions are a family of special functions that are solutions to the modified Bessel's equation:
\[x^2y''+xy'-(x^2+\nu^2)y=0,\]
where $x\in\bbR$ is the independent variable and $\nu\in\bbR$ is called order.
\bit
\item Modified Bessel function of the first kind:
    \[I_{\nu}(x)=\sum_{m=0}^{\infty}\lim_{\mu\to\nu}\frac{1}{m!\Gamma(m+\mu+1)}\qty(\frac{x}{2})^{2m+\nu}.\]
\bit
\item For $\nu\in\bbR\setminus\bbZ$:
    \[I_{\nu}(x)=\sum_{m=0}^{\infty}\frac{1}{m!\Gamma(m+\nu+1)}\qty(\frac{x}{2})^{2m+\nu}.\]
    $I_{\nu}(x)$ and $I_{-\nu}(x)$ are two linearly independent solutions of the modified Bessel's equation. $I_{\nu}(x)$ is defined on $\bbR_{>=0}$ for $\nu\in\bbR_{>0}\setminus\bbZ$, and on $\bbR_{>0}$ for $\nu\in\bbR_{<0}\setminus\bbZ$. $I_{\nu}(0)=0$ for $\nu\in\bbR_{>0}\setminus\bbZ$.
\item For $\nu\in\bbZ$:
    \[I_{\nu}(x)=I_{-\nu}(x)=I_{\nu}(-x)=\sum_{m=0}^{\infty}\lim_{\mu\to\nu}\frac{1}{m!\Gamma(m+\mu+1)}\qty(\frac{x}{2})^{2m+\nu}=\sum_{m=0}^{\infty}\frac{1}{m!(m+\abs{\nu})!}\qty(\frac{x}{2})^{2m+\abs{\nu}}\]
    $I_{\nu}(x)$ is defined on $\bbR$. $I_0(0)=1$. $I_{\nu}(0)=0$ for $\nu\in\bbZ_{\ne0}$. 
\eit
\item Modified Bessel function of the second kind:
    \[K_{\nu}(x)=\lim_{\mu\to\nu}\frac{\pi}{2}\frac{I_{-\mu}(x)-I_{\mu}(x)}{\sin(\mu\pi)}.\]
\bit
\item For $\nu\in\bbR\setminus\bbZ$:
    \[K_{\nu}(x)=\frac{\pi}{2}\frac{I_{-\nu}(x)-I_{\nu}(x)}{\sin(\nu\pi)},\]
    which is a linear combination of $I_{\nu}(x)$ and $I_{-\nu}(x)$.
\item For $\nu\in\bbZ$:
    \[K_{\nu}(x)=\lim_{\mu\to\nu}K_{\mu}(x),\]
    which is a solution that is linearly independent to $I_{\nu}(x)$. $K_{\nu}(x)$ is defined on $\bbR_{>0}$.
    \[\lim_{x\to0^+}K_{\nu}(x)=\infty.\]
\eit
\eit
The general solution is:
\bit
\item For $\nu\in\bbR$:
    \[y(x)=CI_{\nu}(x)+DK_{\nu}(x).\]
\item For $\nu\in\bbR\setminus\bbZ$:
    \[y(x)=CI_{\nu}(x)+DI_{-\nu}(x).\]
\eit
\bpr
\[x^2y''+xy'-(x^2+\nu^2)y=0\]
\[y''+x^{-1}y'-(1+\nu^2x^{-2})y=0\]
$0$ is a regular singular point.
\[y=\sum_{n=0}^{\infty}a_nx^{n+r}\]
\[y'=\sum_{n=0}^{\infty}a_n(n+r)x^{n+r-1}\]
\[y''=\sum_{n=0}^{\infty}a_n(n+r)(n+r-1)x^{n+r-2}\]
\[\sum_{n=0}^{\infty}a_n(n+r)(n+r-1)x^{n+r}+\sum_{n=0}^{\infty}a_n(n+r)x^{n+r}-\sum_{n=0}^{\infty}a_nx^{n+r+2}-\nu^2\sum_{n=0}^{\infty}a_nx^{n+r}=0\]
\[L=a_0r(r-1)x^r+a_0rx^r-a_0x^{r+2}-\nu^2a_0x^r\]
\[I(r)=r^2-\nu^2=0\]
\[r=\pm\nu\]
\[y_1(x)=\sum_{n=0}^{\infty}A_nx^{n+\abs{\nu}}\]
\[y_1'(x)=\sum_{n=0}^{\infty}A_n(n+\abs{\nu})x^{n+\abs{\nu}-1}\]
\[y_1''(x)=\sum_{n=0}^{\infty}A_n(n+\abs{\nu})(n+\abs{\nu}-1)x^{n+\abs{\nu}-2}\]
\[y_2(x)=\sum_{n=0}^{\infty}B_nx^{n-\abs{\nu}}\]
\[y_2'(x)=\sum_{n=0}^{\infty}B_n(n-\abs{\nu})x^{n-\abs{\nu}-1}\]
\[y_2''(x)=\sum_{n=0}^{\infty}B_n(n-\abs{\nu})(n-\abs{\nu}-1)x^{n-\abs{\nu}-2}\]
\[\sum_{n=0}^{\infty}A_n(n+\abs{\nu})(n+\abs{\nu}-1)x^{n+\abs{\nu}}+\sum_{n=0}^{\infty}A_n(n+\abs{\nu})x^{n+\abs{\nu}}-\sum_{n=0}^{\infty}A_nx^{n+\abs{\nu}+2}-\abs{\nu}^2\sum_{n=0}^{\infty}A_nx^{n+\abs{\nu}}=0\]
\[A_1(1+2\abs{\nu})x^{1+\abs{\nu}}+\sum_{n=2}^{\infty}\qty(A_n(n^2+2n\abs{\nu})-A_{n-2})x^{n+\abs{\nu}}=0\]
\[A_n=0,\quad\frac{n+1}{2}\in\bbN\]
\[A_n(n^2+2n\abs{\nu})-A_{n-2}=0,\quad\frac{n}{2}\in\bbN\]
\[A_{2m}=\frac{A_0}{\prod_{i=1}^m(2i)(2i+2\abs{\nu})}=\lim_{\mu\to\abs{\nu}}\frac{A_0\Gamma(\mu+1)}{2^{2m}m!\Gamma(m+\mu+1)},\quad m\in\bbN\]
Let
\[A_0=\lim_{\mu\to\abs{\nu}}\frac{1}{2^\mu\Gamma(\mu+1)}\]
\[A_{2m}=\lim_{\mu\to\abs{\nu}}\frac{1}{2^{2m+\mu}m!\Gamma(m+\mu+1)},\quad m\in\bbN\]
\[y_1(x)=I_{\abs{\nu}}(x)=\sum_{m=0}^{\infty}\lim_{\mu\to\abs{\nu}}\frac{1}{2^{2m+\mu}m!\Gamma(m+\mu+1)}x^{2m+\abs{\nu}}=\sum_{m=0}^{\infty}\lim_{\mu\to\abs{\nu}}\frac{1}{m!\Gamma(m+\mu+1)}\qty(\frac{x}{2})^{2m+\abs{\nu}}\]
\[y_2(x)=I_{-\abs{\nu}}(x)=\sum_{m=0}^{\infty}\lim_{\mu\to\abs{\nu}}\frac{1}{m!\Gamma(m-\mu+1)}\qty(\frac{x}{2})^{2m-\abs{\nu}}\]
For $\nu\in\bbR\setminus\bbZ$, that is,
\[I_{\nu}(x)=\sum_{m=0}^{\infty}\frac{1}{m!\Gamma(m+\nu+1)}\qty(\frac{x}{2})^{2m+\nu}\]
\[I_{-\nu}(x)=\sum_{m=0}^{\infty}\frac{1}{m!\Gamma(m-\nu+1)}\qty(\frac{x}{2})^{2m-\nu}\]
For $2\nu\in\bbZ$, test whether they are are linearly dependent.
\[I_{\nu}'(x)=\sum_{m=0}^{\infty}\lim_{\mu\to\nu}\frac{1}{m!\Gamma(m+\mu+1)}\qty(\frac{x}{2})^{2m+\nu-1}\frac{2m+\nu}{2}\]
\[I_{-\nu}'(x)=\sum_{m=0}^{\infty}\lim_{\mu\to\nu}\frac{1}{m!\Gamma(m-\mu+1)}\qty(\frac{x}{2})^{2m-\nu-1}\frac{2m-\nu}{2}\]
\[W(I_{\nu},I_{-\nu})(x)=\sum_{m=0}^{\infty}\sum_{n=0}^{\infty}\lim_{\mu\to\nu}\frac{(m-n-\nu)}{m!n!\Gamma(m-\mu+1)\Gamma(n+\mu+1)}\qty(\frac{x}{2})^{2m+2n-1}\]
The coefficient of the $x^{-1}$ term is
\[\lim_{\mu\to\nu}\frac{-2\nu}{\Gamma(1-\mu)\Gamma(1+\mu)}=\lim_{\mu\to\nu}\frac{-2}{\Gamma(\mu)\Gamma(1-\mu)}=\frac{-2\sin(\pi\nu)}{\pi}\]
Thus for $\nu\in\bbZ$, $W(I_{\nu},I_{-\nu})=0$,
\[I_{\abs{\nu}}(x)=\sum_{m=0}^{\infty}\frac{1}{m!(m+\abs{\nu})!}\qty(\frac{x}{2})^{2m+\abs{\nu}}\]
and observe that for $k\in\bbN_0$,
\[\lim_{\mu\to\nu}\frac{1}{\Gamma(-k)}=0,\]
thus
\[I_{-\abs{\nu}}(x)=\sum_{m=0}^{\infty}\frac{1}{m!(m+\abs{\nu})!}\qty(\frac{x}{2})^{2m+\abs{\nu}},\]
and
\[I_{-\abs{\nu}}(x)=I_{\abs{\nu}}(x),\]
that is,
\[I_{\nu}(x)=I_{-\nu}(x)=\sum_{m=0}^{\infty}\frac{1}{m!(m+\abs{\nu})!}\qty(\frac{x}{2})^{2m+\abs{\nu}},\]
and a new $y_2(x)$ is
\[y_2(x)=\ln|x|I_{\abs{\nu}}(x)+\sum_{n=0}^{\infty}B_nx^{n-\abs{\nu}}\]
\[y_2'(x)=\ln|x|I_{\abs{\nu}}'(x)+x^{-1}I_{\abs{\nu}}(x)+\sum_{n=0}^{\infty}B_n(n-\abs{\nu})x^{n-\abs{\nu}-1}\]
\[y_2''(x)=\ln|x|I_{\abs{\nu}}''(x)+2x^{-1}I_{\abs{\nu}}'(x)-x^{-2}I_{\abs{\nu}}(x)+\sum_{n=0}^{\infty}B_n(n-\abs{\nu})(n-\abs{\nu}-1)x^{n-\abs{\nu}-2}\]
\[x^2\ln|x|I_{\abs{\nu}}''(x)+2xI_{\abs{\nu}}'(x)-I_{\abs{\nu}}(x)+\sum_{n=0}^{\infty}B_n(n-\abs{\nu})(n-\abs{\nu}-1)x^{n-\abs{\nu}}+x\ln|x|I_{\abs{\nu}}'(x)+I_{\abs{\nu}}(x)+\sum_{n=0}^{\infty}B_n(n-\abs{\nu})x^{n-\abs{\nu}}-x^2\ln|x|I_{\abs{\nu}}(x)-\sum_{n=0}^{\infty}B_nx^{n-\abs{\nu}+2}-\abs{\nu}^2\ln|x|I_{\abs{\nu}}(x)-\abs{\nu}^2\sum_{n=0}^{\infty}B_nx^{n-\abs{\nu}}=0\]
\[2xI_{\abs{\nu}}'(x)+B_1x^{1-\abs{\nu}}+\sum_{n=2}^{\infty}\qty(B_n(n^2-2n\abs{\nu})-B_{n-2})x^{n-\abs{\nu}}=0\]
\[I_{\abs{\nu}}'(x)=\sum_{m=0}^{\infty}\frac{(2m+\abs{\nu})}{2m!(m+\abs{\nu})!}\qty(\frac{x}{2})^{2m+\abs{\nu}-1}\]
\[2xI_{\abs{\nu}}'(x)=\sum_{m=0}^{\infty}\frac{2(2m+\abs{\nu})}{m!(m+\abs{\nu})!}\qty(\frac{x}{2})^{2m+\abs{\nu}}\]
\[\sum_{m=0}^{\infty}\frac{2(2m+\abs{\nu})}{m!(m+\abs{\nu})!}\qty(\frac{x}{2})^{2m+\abs{\nu}}+B_1x^{1-\abs{\nu}}+\sum_{n=2}^{\infty}\qty(B_n(n^2-2n\abs{\nu})-B_{n-2})x^{n-\abs{\nu}}=0\]
TODO
How this = K
\epr
\sssc{Spherical Bessel functions}
Spherical Bessel functions are a family of special functions that are solutions to the equation:
\[x^2y''+2xy'+(x^2-n(n-1))y=0,\]
where $x\in\bbR$ is the independent variable and $n\in\bbZ$.
\bit
\item Spherical Bessel function of the first kind:
    \[j_n(x)=\sqrt{\frac{\pi}{2x}}J_{n+\frac{1}{2}}(x).\]
\item Spherical Bessel function of the second kind:
    \[y_n(x)=\sqrt{\frac{\pi}{2x}}Y_{n+\frac{1}{2}}(x)=(-1)^{n+1}\sqrt{\frac{\pi}{2x}}J_{-n-\frac{1}{2}}(x),\]
    which is linearly independent from $j_n(x)$.
\eit
\[j_n(x)=(-1)^ny_{-n-1}(x)\]
\[y_n(x)=(-1)^{n+1}j_{-n-1}(x)\]
\[j_0(x)=\frac{\sin(x)}{x}\]
\[y_0(x)=-\frac{\cos(x)}{x}\]
The general solution is:
\[y=Cj_n(x)+Dy_n(x).\]
\bpr
\[u=x^{0.5}y\]
\[u'=x^{0.5}y'+0.5x^{-0.5}y\]
\[u''=x^{0.5}y''+x^{-0.5}y'-0.25x^{-1.5}y\]
\[x^{1.5}u''=x^2y''+xy'-0.25y\]
\[x^{0.5}u'=xy'+0.5y\]
\[x^{1.5}u''+x^{0.5}u'+(x^2-n(n-1)-0.25)x^{-0.5}u=0,\]
\[x^2u''+xu'+(x^2-(n-0.5)^2)u=0\]
\[u=CJ_{n+0.5}(x)+DY_{n+0.5}(x)\]
\epr
\sssc{Parametric Bessel equation}
ODE:
\[x^2y''+xy'+(\alpha^2x^2-\nu^2)y=0,\quad\alpha\ne 0\]
Solution:
\[y(x)=CJ_{\nu}(\alpha x)+DY_{\nu}(\alpha x).\]
\bpr
Let
\[z=\alpha x\]
\[y(x)=u(z)\]
\[y'(x)=\alpha u'(z)\]
\[y''(x)=\alpha^2u''(z)\]
\[z^2u''+zu'+(z^2-\nu^2)u=0\]
\[u(z)=CJ_{\nu}(z)+DY_{\nu}(z).\]
\epr
\sssc{Parametric Modified Bessel equation}
ODE:
\[x^2y''+xy'-(\alpha^2x^2+\nu^2)y=0,\quad\alpha\ne 0\]
Solution:
\[y(x)=CI_{\nu}(\alpha x)+DK_{\nu}(\alpha x).\]
\bpr
Let
\[z=\alpha x\]
\[y(x)=u(z)\]
\[y'(x)=\alpha u'(z)\]
\[y''(x)=\alpha^2u''(z)\]
\[z^2u''+zu'-(z^2+\nu^2)u=0\]
\[u(z)=CI_{\nu}(z)+DK_{\nu}(z).\]
\epr
\sssc{Derivative of Bessel function of first kind}
\[\dv{}{x}\qty(x^{-\nu}J_{\nu}(x))=-x^{-\nu}J_{\nu+1}(x).\]
\bpr
\[\ba
\dv{}{x}\qty(x^{-\nu}J_{\nu}(x))&=-\nu x^{-\nu-1}J_{\nu}(x)+x^{-\nu}J_{\nu}'(x)\\
&=-\nu x^{-\nu-1}\sum_{m=0}^{\infty}\lim_{\mu\to\nu}\frac{(-1)^m}{m!\Gamma(m+\mu+1)}\qty(\frac{x}{2})^{2m+\nu}+x^{-\nu}\sum_{m=0}^{\infty}\lim_{\mu\to\nu}\frac{(-1)^m(2m+\nu)}{2m!\Gamma(m+\mu+1)}\qty(\frac{x}{2})^{2m+\nu-1}\\
&=-x^{-\nu}\sum_{m=0}^{\infty}\lim_{\mu\to\nu}\frac{(-1)^m\nu}{2m!\Gamma(m+\mu+1)}\qty(\frac{x}{2})^{2m+\nu-1}+x^{-\nu}\sum_{m=0}^{\infty}\lim_{\mu\to\nu}\frac{(-1)^m(2m+\nu)}{2m!\Gamma(m+\mu+1)}\qty(\frac{x}{2})^{2m+\nu-1}\\
&=x^{-\nu}\sum_{m=1}^{\infty}\lim_{\mu\to\nu}\frac{(-1)^m}{(m-1)!\Gamma(m+\mu+1)}\qty(\frac{x}{2})^{2m+\nu-1}\\
&=x^{-\nu}\sum_{m=0}^{\infty}\lim_{\mu\to\nu}\frac{(-1)^{m+1}}{m!\Gamma(m+\mu+2)}\qty(\frac{x}{2})^{2m+\nu+1}\\
&=-x^{-\nu}J_{\nu+1}(x)
\ea\]
\epr
\[\dv{}{x}\qty(x^{\nu}J_{\nu}(x))=x^{\nu}J_{\nu-1}(x).\]
\bpr
\[\ba
\dv{}{x}\qty(x^{\nu}J_{\nu}(x))&=\nu x^{\nu-1}J_{\nu}(x)+x^{\nu}J_{\nu}'(x)\\
&=\nu x^{\nu-1}\sum_{m=0}^{\infty}\lim_{\mu\to\nu}\frac{(-1)^m}{m!\Gamma(m+\mu+1)}\qty(\frac{x}{2})^{2m+\nu}+x^{\nu}\sum_{m=0}^{\infty}\lim_{\mu\to\nu}\frac{(-1)^m(2m+\nu)}{2m!\Gamma(m+\mu+1)}\qty(\frac{x}{2})^{2m+\nu-1}\\
&=x^{\nu}\sum_{m=0}^{\infty}\lim_{\mu\to\nu}\frac{(-1)^m\nu}{2m!\Gamma(m+\mu+1)}\qty(\frac{x}{2})^{2m+\nu-1}+x^{\nu}\sum_{m=0}^{\infty}\lim_{\mu\to\nu}\frac{(-1)^m(2m+\nu)}{2m!\Gamma(m+\mu+1)}\qty(\frac{x}{2})^{2m+\nu-1}\\
&=x^{\nu}\sum_{m=0}^{\infty}\lim_{\mu\to\nu}\frac{(-1)^m}{m!\Gamma(m+\mu)}\qty(\frac{x}{2})^{2m+\nu-1}\\
&=x^{\nu}J_{\nu-1}(x)
\ea\]
\epr
\[xJ_{\nu}'(x)=\nu J_{\nu}(x)-xJ_{\nu+1}(x).\]
\bpr
\[\ba
xJ_{\nu}'(x)&=\sum_{m=0}^{\infty}\lim_{\mu\to\nu}\frac{(-1)^m(2m+\nu)}{m!\Gamma(m+\mu+1)}\qty(\frac{x}{2})^{2m+\nu}\\
&=\nu\sum_{m=0}^{\infty}\lim_{\mu\to\nu}\frac{(-1)^m}{m!\Gamma(m+\mu+1)}\qty(\frac{x}{2})^{2m+\nu}+2\sum_{m=1}^{\infty}\lim_{\mu\to\nu}\frac{(-1)^m}{(m-1)!\Gamma(m+\mu+1)}\qty(\frac{x}{2})^{2m+\nu}\\
&=\nu J_{\nu}(x)+x\sum_{m=1}^{\infty}\lim_{\mu\to\nu}\frac{(-1)^m}{(m-1)!\Gamma(m+\mu+1)}\qty(\frac{x}{2})^{2m+\nu-1}\\
&=\nu J_{\nu}(x)+x\sum_{m=0}^{\infty}\lim_{\mu\to\nu}\frac{(-1)^{m+1}}{m!\Gamma(m+\mu+2)}\qty(\frac{x}{2})^{2m+\nu+1}\\
&=\nu J_{\nu}(x)-xJ_{\nu+1}(x)
\ea\]
\epr
\[xJ_{\nu}'(x)=xJ_{\nu-1}(x)-\nu J_{\nu}(x).\]
\bpr
\[\ba
xJ_{\nu}'(x)&=\sum_{m=0}^{\infty}\lim_{\mu\to\nu}\frac{(-1)^m(2m+\nu)}{m!\Gamma(m+\mu+1)}\qty(\frac{x}{2})^{2m+\nu}\\
&=2\sum_{m=0}^{\infty}\lim_{\mu\to\nu}\frac{(-1)^m(m+\nu)}{m!\Gamma(m+\mu+1)}\qty(\frac{x}{2})^{2m+\nu}-\nu\sum_{m=0}^{\infty}\lim_{\mu\to\nu}\frac{(-1)^m}{m!\Gamma(m+\mu+1)}\qty(\frac{x}{2})^{2m+\nu}\\
&=x\sum_{m=0}^{\infty}\lim_{\mu\to\nu}\frac{(-1)^m}{m!\Gamma(m+\mu)}\qty(\frac{x}{2})^{2m+\nu-1}-\nu\sum_{m=0}^{\infty}\lim_{\mu\to\nu}\frac{(-1)^m}{m!\Gamma(m+\mu+1)}\qty(\frac{x}{2})^{2m+\nu}\\
&=xJ_{\nu-1}(x)-\nu J_{\nu}(x)
\ea\]
\epr
\[J_{\nu}'(x)=\frac{1}{2}\qty(J_{\nu-1}(x)-J_{\nu+1}(x)).\]
\bpr
\[\ba
J_{\nu}'(x)&=\frac{1}{2}\sum_{m=0}^{\infty}\lim_{\mu\to\nu}\frac{(-1)^m(2m+\nu)}{m!\Gamma(m+\mu+1)}\qty(\frac{x}{2})^{2m+\nu-1}\\
&=\frac{1}{2}\qty(\sum_{m=0}^{\infty}\lim_{\mu\to\nu}\frac{(-1)^m(m+\nu)}{m!\Gamma(m+\mu+1)}\qty(\frac{x}{2})^{2m+\nu-1}+\sum_{m=0}^{\infty}\lim_{\mu\to\nu}\frac{(-1)^mm}{m!\Gamma(m+\mu+1)}\qty(\frac{x}{2})^{2m+\nu-1})\\
&=\frac{1}{2}\qty(\sum_{m=0}^{\infty}\lim_{\mu\to\nu}\frac{(-1)^m}{m!\Gamma(m+\mu)}\qty(\frac{x}{2})^{2m+\nu-1}+\sum_{m=1}^{\infty}\lim_{\mu\to\nu}\frac{(-1)^{m+1}}{m!\Gamma(m+\mu+2)}\qty(\frac{x}{2})^{2m+\nu+1})\\
&=\frac{1}{2}\qty(J_{\nu-1}(x)-J_{\nu+1}(x))
\ea\]
\epr
\sssc{Derivative of modified Bessel function of first kind}
\[\dv{}{x}\qty(x^{-\nu}I_{\nu}(x))=x^{-\nu}I_{\nu+1}(x).\]
\bpr
\[\ba
\dv{}{x}\qty(x^{-\nu}I_{\nu}(x))&=-\nu x^{-\nu-1}I_{\nu}(x)+x^{-\nu}I_{\nu}'(x)\\
&=-\nu x^{-\nu-1}\sum_{m=0}^{\infty}\lim_{\mu\to\nu}\frac{1}{m!\Gamma(m+\mu+1)}\qty(\frac{x}{2})^{2m+\nu}+x^{-\nu}\sum_{m=0}^{\infty}\lim_{\mu\to\nu}\frac{(2m+\nu)}{2m!\Gamma(m+\mu+1)}\qty(\frac{x}{2})^{2m+\nu-1}\\
&=-x^{-\nu}\sum_{m=0}^{\infty}\lim_{\mu\to\nu}\frac{\nu}{2m!\Gamma(m+\mu+1)}\qty(\frac{x}{2})^{2m+\nu-1}+x^{-\nu}\sum_{m=0}^{\infty}\lim_{\mu\to\nu}\frac{(2m+\nu)}{2m!\Gamma(m+\mu+1)}\qty(\frac{x}{2})^{2m+\nu-1}\\
&=x^{-\nu}\sum_{m=1}^{\infty}\lim_{\mu\to\nu}\frac{1}{(m-1)!\Gamma(m+\mu+1)}\qty(\frac{x}{2})^{2m+\nu-1}\\
&=x^{-\nu}\sum_{m=0}^{\infty}\lim_{\mu\to\nu}\frac{1}{m!\Gamma(m+\mu+2)}\qty(\frac{x}{2})^{2m+\nu+1}\\
&=x^{-\nu}I_{\nu+1}(x)
\ea\]
\epr
\[\dv{}{x}\qty(x^{\nu}I_{\nu}(x))=x^{\nu}I_{\nu-1}(x).\]
\bpr
\[\ba
\dv{}{x}\qty(x^{\nu}I_{\nu}(x))&=\nu x^{\nu-1}I_{\nu}(x)+x^{\nu}I_{\nu}'(x)\\
&=\nu x^{\nu-1}\sum_{m=0}^{\infty}\lim_{\mu\to\nu}\frac{1}{m!\Gamma(m+\mu+1)}\qty(\frac{x}{2})^{2m+\nu}+x^{\nu}\sum_{m=0}^{\infty}\lim_{\mu\to\nu}\frac{(2m+\nu)}{2m!\Gamma(m+\mu+1)}\qty(\frac{x}{2})^{2m+\nu-1}\\
&=x^{\nu}\sum_{m=0}^{\infty}\lim_{\mu\to\nu}\frac{\nu}{2m!\Gamma(m+\mu+1)}\qty(\frac{x}{2})^{2m+\nu-1}+x^{\nu}\sum_{m=0}^{\infty}\lim_{\mu\to\nu}\frac{(2m+\nu)}{2m!\Gamma(m+\mu+1)}\qty(\frac{x}{2})^{2m+\nu-1}\\
&=x^{\nu}\sum_{m=0}^{\infty}\lim_{\mu\to\nu}\frac{1}{m!\Gamma(m+\mu)}\qty(\frac{x}{2})^{2m+\nu-1}\\
&=x^{\nu}I_{\nu-1}(x)
\ea\]
\epr
\[xI_{\nu}'(x)=\nu I_{\nu}(x)+xI_{\nu+1}(x).\]
\bpr
\[\ba
xI_{\nu}'(x)&=\sum_{m=0}^{\infty}\lim_{\mu\to\nu}\frac{(2m+\nu)}{m!\Gamma(m+\mu+1)}\qty(\frac{x}{2})^{2m+\nu}\\
&=\nu\sum_{m=0}^{\infty}\lim_{\mu\to\nu}\frac{1}{m!\Gamma(m+\mu+1)}\qty(\frac{x}{2})^{2m+\nu}+2\sum_{m=1}^{\infty}\lim_{\mu\to\nu}\frac{1}{(m-1)!\Gamma(m+\mu+1)}\qty(\frac{x}{2})^{2m+\nu}\\
&=\nu I_{\nu}(x)+x\sum_{m=1}^{\infty}\lim_{\mu\to\nu}\frac{1}{(m-1)!\Gamma(m+\mu+1)}\qty(\frac{x}{2})^{2m+\nu-1}\\
&=\nu I_{\nu}(x)+x\sum_{m=0}^{\infty}\lim_{\mu\to\nu}\frac{1}{m!\Gamma(m+\mu+2)}\qty(\frac{x}{2})^{2m+\nu+1}\\
&=\nu I_{\nu}(x)+xI_{\nu+1}(x)
\ea\]
\epr
\[xI_{\nu}'(x)=xI_{\nu-1}(x)-\nu I_{\nu}(x).\]
\bpr
\[\ba
xI_{\nu}'(x)&=\sum_{m=0}^{\infty}\lim_{\mu\to\nu}\frac{(2m+\nu)}{m!\Gamma(m+\mu+1)}\qty(\frac{x}{2})^{2m+\nu}\\
&=2\sum_{m=0}^{\infty}\lim_{\mu\to\nu}\frac{(m+\nu)}{m!\Gamma(m+\mu+1)}\qty(\frac{x}{2})^{2m+\nu}-\nu\sum_{m=0}^{\infty}\lim_{\mu\to\nu}\frac{1}{m!\Gamma(m+\mu+1)}\qty(\frac{x}{2})^{2m+\nu}\\
&=x\sum_{m=0}^{\infty}\lim_{\mu\to\nu}\frac{1}{m!\Gamma(m+\mu)}\qty(\frac{x}{2})^{2m+\nu-1}-\nu\sum_{m=0}^{\infty}\lim_{\mu\to\nu}\frac{1}{m!\Gamma(m+\mu+1)}\qty(\frac{x}{2})^{2m+\nu}\\
&=xI_{\nu-1}(x)-\nu I_{\nu}(x)
\ea\]
\epr
\[I_{\nu}'(x)=\frac{1}{2}\qty(I_{\nu-1}(x)+I_{\nu+1}(x)).\]
\bpr
\[\ba
I_{\nu}'(x)&=\frac{1}{2}\sum_{m=0}^{\infty}\lim_{\mu\to\nu}\frac{(2m+\nu)}{m!\Gamma(m+\mu+1)}\qty(\frac{x}{2})^{2m+\nu-1}\\
&=\frac{1}{2}\qty(\sum_{m=0}^{\infty}\lim_{\mu\to\nu}\frac{(m+\nu)}{m!\Gamma(m+\mu+1)}\qty(\frac{x}{2})^{2m+\nu-1}+\sum_{m=0}^{\infty}\lim_{\mu\to\nu}\frac{m}{m!\Gamma(m+\mu+1)}\qty(\frac{x}{2})^{2m+\nu-1})\\
&=\frac{1}{2}\qty(\sum_{m=0}^{\infty}\lim_{\mu\to\nu}\frac{1}{m!\Gamma(m+\mu)}\qty(\frac{x}{2})^{2m+\nu-1}+\sum_{m=1}^{\infty}\lim_{\mu\to\nu}\frac{1}{m!\Gamma(m+\mu+2)}\qty(\frac{x}{2})^{2m+\nu+1})\\
&=\frac{1}{2}\qty(I_{\nu-1}(x)+I_{\nu+1}(x))
\ea\]
\epr
\sssc{Legendre polynomial}
Legendre polynomials are the polynomial solutions of Legendre's differential equations for $\abs{x}<1$:
\[(1-x^2)P_n''(x)-2xP_n'(x)+n(n+1)P_n(x)=0,\]
where $n\in\bbN_0$.
\[P_n(x)=\frac{1}{2^n}\sum_{k=0}^{\left\lfloor\frac{n}{2}\right\rfloor}(-1)^k\binom{2n-2k}{n}\binom{n}{k}x^{n-2k}\]
\begin{proof}
\[P_n''(x)-\frac{2x}{1-x^2}P_n'(x)+\frac{n(n+1)}{1-x^2}P_n(x)=0.\]
\(\pm 1\) are two regular singular points.

Use power series method at $0$.
\[y(x)=\sum_{m=0}^{\infty}a_mx^m\]
\[y'(x)=\sum_{m=1}^{\infty}a_mmx^{m-1}\]
\[y''(x)=\sum_{m=2}^{\infty}a_mm(m-1)x^{m-2}\]
\[\sum_{m=2}^{\infty}a_mm(m-1)x^{m-2}-\sum_{m=2}^{\infty}a_mm(m-1)x^m-2\sum_{m=1}^{\infty}a_mmx^m+n(n+1)\sum_{m=0}^{\infty}a_mx^m=0\]
\[2a_2+n(n+1)a_0+(6a_3-2a_1+n(n+1)a_1)x+\sum_{m=2}^{\infty}(a_{m+2}(m+2)(m+1)-a_m(m+1)+n(n+1)a_m)x^m=0\]
\[a_2=-\frac{n(n+1)}{2}a_0\]
\[a_3=-\frac{(n-1)(n+2)}{6}a_1\]
\[a_{m+2}=-\frac{(n-m)(n+m+1)}{(m+1)(m+2)}a_m,\quad m\in\bbN_0\]
\[a_m=-\frac{(m+1)(m+2)}{(n-m)(n+m+1)}a_{m+2},\quad m\in\bbN_0\]
$a_m$ terminates at $m=n$, which yields the polynomial solution.

For even $n$, let $c=\frac{n}{2}$.
\[a_n=(-1)^c\frac{(2n)!(c!)^2}{(n!)^3}a_0\]
Take
\[a_0=(-1)^c\frac{n!}{2^n(c!)^2}\]
\[a_n=\frac{(2n)!}{2^n(n!)^2}\]
\[a_{n-2k}=(-1)^k\frac{n!}{(n-2k)!}\frac{1}{2^kk!}\frac{2^nn!(2n-2k)!}{2^{n-k}(n-k)!(2n)!}a_n\]
\[a_{n-2k}=(-1)^k\frac{(2n-2k)!}{(n-2k)!}\frac{1}{k!(n-k)!}\frac{1}{2^n}=\frac{(-1)^k}{2^n}\binom{2n-2k}{n}\binom{n}{k}\]
For odd $n$, let $c=\frac{n-1}{2}$.
\[a_n=(-1)^c\frac{(2n)!(c!)^2}{2(n!)^3}a_1\]
Take
\[a_1=(-1)^c\frac{n!}{2^{n-1}(c!)^2}\]
\[a_n=\frac{(2n)!}{2^n(n!)^2}\]
\[a_{n-2k}=(-1)^k\frac{n!}{(n-2k)!}\frac{1}{2^kk!}\frac{2^nn!(2n-2k)!}{2^{n-k}(n-k)!(2n)!}a_n\]
\[a_{n-2k}=(-1)^k\frac{(2n-2k)!}{(n-2k)!}\frac{1}{k!(n-k)!}\frac{1}{2^n}=\frac{(-1)^k}{2^n}\binom{2n-2k}{n}\binom{n}{k}\]
Thus, for all $n\in\bbN_0$,
\[P_n(x)=\frac{1}{2^n}\sum_{k=0}^{\left\lfloor\frac{n}{2}\right\rfloor}(-1)^k\binom{2n-2k}{n}\binom{n}{k}x^{n-2k}\]
\end{proof}
\sssc{Rodrigues' formula}
The Legendre polynomials are given by
\[P_n(x)=\frac{1}{2^nn!}\dv[n]{x}\qty(x^2-1)^n\]
\begin{proof}
TODO 
\end{proof}
\section{Linear First-order System}
\ssc{Basics}
\sssc{System of linear first-order ODEs or linear first-order system or linear system}
A linear first-order system can be written in matrix form
\[\mb{y}'(x)=\mb{A}(x)\mb{y}(x)+\mb{f}(x),\]
where
\[\mb{y}(x)=\begin{bmatrix}y_1(x)\\\vdots\\y_n(x)\end{bmatrix}\]
is called solution vector,
\[\mb{A}(x)=\begin{bmatrix}a_{1,1}(x) & \cdots & a_{1,n}(x)\\\vdots & \ddots & \vdots\\a_{n,1}(x) & \cdots & a_{n,n}(x)\end{bmatrix}\]
is called coefficient matrix, and
\[\mb{f}(x)=\begin{bmatrix}f_1(x)\\\vdots\\f_n(x)\end{bmatrix}\]
is called forcing term.

And an IC of an IVP of it can be written in the form
\[\mb{y}(x_0)=\mb{y}_0=\begin{bmatrix}\gamma_1\\\vdots\\\gamma_n\end{bmatrix}\]
where $x_0$ is a constant in the domain of $\mb{A}$ and $\mb{f}$, and $\mb{y}_0$ is a constant vector.

A homogeneous one is called homogeneous system.
\sssc{Solving by Elimination}
\ben
\item Write the system in terms of differential operator $D$.
\item Manipulate the polynomials of $D$ and unknowns such that each linear equation contains only each unknown.
\item Solve all of them.
\een
\sssc{Wrońskian}
Suppose $n$ functions $\mb{y}_i\colon\Omega_i\subseteq\bbR\to\bbR^n$ for all $i\in\bbN_{\le n}$ are solution vectors to a homogeneous system on an interval and
\[\mb{y}_i=\begin{bmatrix}y_{1i}\\\vdots\\y_{ni}\end{bmatrix}.\]
Then the Wrońskian of them is defined as
\[W(\mb{y}_1,\ldots,\mb{y}_n)=\begin{vmatrix}\mb{y}_1&\ldots&\mb{y}_n\end{vmatrix}=\begin{vmatrix}y_{11}&\ldots&y_{1n}\\\vdots&\ddots&\vdots\\y_{n1}&\ldots&y_{nn}\end{vmatrix}.\]
\sssc{Non-zero Wrońskian implies linear independence}
For $n$ functions $\mb{y}_i\colon\Omega_i\subseteq\bbR\to\bbR^n$ for all $i\in\bbN_{\le n}$, if $W(\mb{y}_1,\ldots,\mb{y}_n)\ne0$, then they are linearly independent.
\begin{proof}
If they are not linearly independent, there exists not all zero $c_1,\ldots,c_n$ such that
\[\sum_{i=1}^nc_i\mb{y}_i=0,\]
that is,
\[\begin{bmatrix}y_{11}&\ldots&y_{1n}\\\vdots&\ddots&\vdots\\y_{n1}&\ldots&y_{nn}\end{bmatrix}=\begin{bmatrix}c_1\\c_2\\\vdots\\c_n\end{bmatrix}=\begin{bmatrix}0\\0\\\vdots\\0\end{bmatrix}\]
All $c_i=0$ if $W(\mb{y}_1,\ldots,\mb{y}_n)\ne 0$ by Cramer's rule, thus $W(\mb{y}_1,\ldots,\mb{y}_n)=0$.
\end{proof}
\sssc{Homogeneity}
A linear first-order system is called a homogeneous linear first-order system iff the forcing term is 0 for all $x$; otherwise, it is called a non-homogeneous linear first-order system.
\sssc{Complementary or associated (homogeneous) function, natural, complementary, or homogeneous solution or function, and particularly solution}
Given an $n$th-order linear first-order system 
\[\mb{y}'(x)=\mb{A}(x)\mb{y}(x)+\mb{f}(x),\]
\[\mb{y}'(x)=\mb{A}(x)\mb{y}(x),\]
is called its complementary or associated (homogeneous) function, and the general solutions of the original equation is the sum of the general solutions $\mb{y}_c$ of the complementary function, called natural, complementary, or homogeneous solution or function, and a particular solution $\mb{y}_p$.
\sssc{Superposition principle}
If $\mb{y}_1(x)$ and $\mb{y}_2(x)$ are both solutions to a linear first-order system on an interval, then $c_1\mb{y}_1(x)+c_2\mb{y}_2(x)$ is also a solution to it on the interval for any $c_1,c_2\in\bbR$.
\ssc{Homogeneous linear first-order system}
TODO
\ssc{Solving linear ODE by conversion to linear first-order system}
\sssc{Converting linear ODE to linear first-order system}
\[\sum_{i=0}^na_i(x)y^{(i)}=f(x),\]
TODO
\sssc{Existence and uniqueness theorem for IVP of linear ODE}
TODO, see GPT
Let $n\in\bbZ$, $a_i(x)$ for all $i\in[0,n]\cap\bbZ$ and $g(x)$ be continuous on an interval $I$, $a_n(x)\ne0$ on $I$, and $x_0\in I$. Then the IVP of
\[\sum_{i=0}^na_i(x)\dv[i](y)(x)=g(x)\]
subjected to $y^{(i)}(x_0)=y_i$ for all $i\in[0,n-1]\cap\bbZ$ has a solution $y(x)$ on $I$ and the solution is unique.
existence and uniqueness theorem's hypothesis that the leading coefficient is nonzero at $x_0$
\sssc{BVP of linear ODE}
TODO
\section{Dynamical system (DS) (動態系統 or 動力系統)}
\ssc{Dynamical system (DS)}
\sssc{Dynamical system}
A dynamical system, or simply system, is a tuple $(T, X, \Phi)$ where $T$ is a monoid, written additively and with the independent variable $t$ in it called evolution parameter (演化參數) or time, $X$ is a non-empty set, called the phase space (相空間) or state space (狀態空間) with the independent variable $x$ in it called (system) phase (相) or state (狀態), and $\Phi$ is a function, called the evolution function (演化函數) or flow (map),
\[\Phi(t,x)\colon U\subseteq(T\times X)\to X\]
where $U$ is such that the second projection map $\pi_2$ for $T\times X$ satisfies
\[\pi_2(U)=X,\]
and for any $x \in X$, $\Phi (0,x)=x$.
\sssc{Time-invariant or autonomous dynamical system}
A dynamic system $(T, X, \Phi)$ is called time-invariant or autonomous iff it satisfies the semigroup property, translation property, shift property, time-invariance, or autonomy that
\[\Phi (t_{2},\Phi (t_{1},x))=\Phi (t_{2}+t_{1},x)\]
for any $t_{1},t_{2}+t_{1}\in I(x)$ and $t_{2}\in I(\Phi (t_{1},x))$, in which
\[I(x)\coloneq\{t\in T\mid(t,x)\in U\}\]
for any $x\in X$.

Otherwise, it is called time-variant or non-autonomous.
\sssc{Orbit (軌道)}
Given a dynamical system $(T, X, \Phi)$ with
\[\Phi(t,x)\colon U\subseteq(T\times X)\to X\]
and
\[I(x)\coloneq\{t\in T\mid(t,x)\in U\}.\]
The set
\[\gamma(x_0)\coloneq\{\Phi (t,x)\colon t\in I(x)\}\subset X\]
is called the orbit through $x$.

An orbit which consists of a single point is called constant orbit. A non-constant orbit is called closed or periodic if there exists $t\neq 0\in I(x)$ such that $\Phi (t,x)=x$.
\sssc{(Phase) trajectory (function)}
Given a dynamical system $(T, X, \Phi)$, for any $x_0\in X$, the (phase) trajectory (function) starting at $x_0$ is a function
\[x\colon t\to\gamma(x_0);\;t\mapsto\Phi(t,x_0).\]
\sssc{Phase portrait (相圖)}
A phase portrait is a geometric representation of a collection of phase trajectories of an autonomous dynamical system, typically all phase trajectories starting at any $x\in X$.
\sssc{Attractor (吸引子)}
For a dynamical system $(\mathbb{R}, X, \Phi)$ with $D\subseteq X$ and $\Phi(t,x)\colon \mathbb{R}_{\geq 0}\times D\to X$, a non-empty subset $A$ of $D$ is called an attractor if and only if it satisfies the following three conditions:
\bit
\item Forward invariance:
    \[\forall a\in A\colon\forall t>0\colon\Phi(t,a)\in A.\]
\item Attraction: There exists a neighborhood of $A$, called the basin of attraction for $A$ and denoted $B(A)$, such that for any $b\in B(A)$ and any open neighborhood $N$ of $A$, there is a positive constant $T$ such that $\forall t>T\colon\Phi(t,b)\in N$.
\item Minimality: There is no non-empty proper subset of $A$ that satisfies the above two conditions.
\eit
A point $a$ in $D$ is called an attractor iff $\{a\}$ is an attractor.
\sssc{Repeller}
For a dynamical system $(\mathbb{R}, X, \Phi)$ with $D\subseteq X$ and $\Phi(t,x)\colon \mathbb{R}_{\geq 0}\times D\to X$, a non-empty subset $R$ of $X$ is called a repeller if and only if it satisfies the following three conditions:
\bit
\item Forward invariance:
    \[\forall r\in R\colon\forall t>0\colon\Phi(t,r)\in R.\]
\item Repulsion: There exists a neighborhood of $R$, called the basin of repulsion for $R$ and denoted $B(R)$, such that for any $b\in B(R)\setminus R$ and any open neighborhood $N$ of $R$, there is a positive constant $T$ such that $\forall t>T\colon\Phi(t,b)\notin N$.
\item Minimality: There is no non-empty proper subset of $R$ that satisfies the above two conditions.
\eit
A point $r$ in $D$ is called a repeller iff $\{r\}$ is a repeller.
\sssc{Continuous-time dynamical system}
A continuous-time dynamical system is a dynamical system $(\mathbb{R}, X, \Phi)$ with $X$ being a manifold locally diffeomorphic to a Banach space and $\Phi(t,x)\colon A\times X\to X$ with $\mathbb{R}_{\geq 0}\subseteq A\subseteq\mathbb{R}$ such that $\pdv{\Phi}{t}\big\vert_{t=0}$ exists for all $x\in X$.
\sssc{Discrete-time dynamical system}
A discrete-time dynamical system is a dynamical system $(\mathbb{Z}, X, \Phi)$ with $X$ being a manifold locally diffeomorphic to a Banach space and $\Phi(t,x)\colon A\times X\to X$ with $\mathbb{N}_0\subseteq A\subseteq\mathbb{Z}$.
\ssc{Continuous-time dynamical systems}
Below, we will discuss continuous-time dynamical systems in Euclidean vector spaces, that are, continuous-time dynamical systems whose phase space is $\bbR^n$ for some $n\in\bbN$.
\sssc{Systems of first-order ODEs}
Consider a system of first-order ODEs:
\[\dv{x}{t}=f(t,x), \quad x\in\mathbb{R}^n\land f\colon\mathbb{R}\times\mathbb{R}^n\to\mathbb{R}^n.\]
Let $U\subseteq\mathbb{R}\times\mathbb{R}^n$ be the set such that for any $(t_0,x_0)$ in $U$, $t_0$ is in the interval of definition of the solution of it with the initial condition $x(0)=x_0$ and that the solution is unique, and $x(t;x_0)$ be the solution of it with the initial condition $x(0)=x_0$.

Then the dynamical system represented by it has phase space $\mathbb{R}^n$, evolution parameter $t$, phase $x$, and evolution function $\Phi\colon U\to\mathbb{R}^n;\Phi(t,x_0)=x(t;x_0)$. The system of ODEs is called the evolution rule of the dynamical system, and the dynamical system is called to be defined by the system of ODEs. Each solution of the system of ODEs is a phase trajectory of the dynamical system.
\ssc{Autonomous continuous-time dynamical systems}
Below, we will discuss autonomous continuous-time dynamical systems in Euclidean vector spaces, that are, autonomous continuous-time dynamical systems whose phase space is $\bbR^n$ for some $n\in\bbN$.
\sssc{Equilibria, fixed points, or fixpoints of autonomous continuous-time dynamical systems}
An equilibrium, fixed point, or fixpoint of an autonomous continuous-time dynamical system $(T, X, \Phi)$ is a point $x^*\in X$ such that the vector field $f(x)=\pdv{\Phi}{t}\big\vert_{t=0}$ vanishes there, that is, $f(x^*)=0$.
\sssc{Autonomous system of ODEs}
An autonomous system of ODEs defines an autonomous continuous-time dynamical system, and the equilibria of the dynamical system are the critical points of the system of ODEs.
\sssc{Phase portrait}
For an autonomous continuous-time dynamical system, one usually plots $X$, marks equilibria, and draw arrows at each point $x$ in a grid indicating $f(x)$ at it. When $X=\mathbb{R}$, a phase potrait is also called a phase line.
\sssc{Linear autonomous continuous-time dynamical system}
An autonomous continuous-time dynamical system $(\mathbb{R}, \mathbb{R}^n, \Phi(t,x))$ is called linear iff its evolution rule is a system of linear ODEs. Let the forcing term be $q(x)=Cx^\top$ where $C$ is a constant matrix in $\mathbb{R}^{n\times n}$, then $C$ is called the dynamics matrix of the system, and the eigenvalues and eigen vectors of $C$ are called the eigenvalues and eigen vectors of the system.
\ssc{Stability of equilibria}
Let $x^*$ be an equilibrium of an autonomous continuous-time dynamical system and $A$ be the domain of the evolution function.
\sssc{(Lyapunov) stable (（李亞普諾夫）穩定)}
$x^*$ is called to be (Lyapunov) stable iff
\[\forall\varepsilon>0\colon\exists\delta>0\text{\ s.t.\ }\|x(0)-x^*\|<\delta\implies\forall t\in A\colon\|x(t)-x^*\|<\varepsilon.\]
\sssc{Asymptotically stable (漸近穩定)}
$x^*$ is called to be asymptotically stable iff
\[\exists\delta>0\text{\ s.t.\ }\|x(0)-x^*\|<\delta\implies\lim_{t\to\infty}x(t)=x^*.\]
An asymptotically stable equilibrium is necessarily (Lyapunov) stable and an attractor.
\sssc{Exponentially stable (指數穩定)}
$x^*$ is called to be exponentially stable iff
\[\exists M>0,\alpha>0,\delta>0\text{\ s.t.\ }\|x(0)-x^*\|<\delta\implies\forall t\in A\|x(t)-x^*\|\leq\|x(0)-x^*\|e^{-\alpha t}.\]
An exponentially stable equilibrium is necessarily asymptotically stable.
\sssc{Semi-stable (半穩定)}
$x^*$ is called to be semi-stable iff it is (Lyapunov) stable and
\[(\exists\delta>0\text{\ s.t.\ }x(0)-x^*<\delta\implies\lim_{t\to\infty}x(t)=x^*)\lor (\exists\delta>0\text{\ s.t.\ }x^*-x(0)<\delta\implies\lim_{t\to\infty}x(t)=x^*).\]
\sssc{Marginally stable (臨界穩定)}
For linear autonomous continuous-time dynamical system, $x^*$ is called to be marginally stable iff it is (Lyapunov) stable but not asymptotically stable.
\sssc{Unstable (不穩定)}
$x^*$ is called to be unstable iff it is not (Lyapunov) stable.

An equilibrium that is in a repeller is necessarily unstable.
\section{Numerical Method}
\ssc{First-order ODE}
\sssc{(Forward) Euler('s) method}
Approximate values for the solution of the initial-value problem $y'=f(x,y)$, $y(x_0)=y_0$, with step size $h$, at $x_n=x_{n-1}+h$, are
\[y_n=y_{n-1}+hf(x_{n-1},y_{n-1}),\quad n\in\bbN.\]
\sssc{Backward Euler('s) method}
Approximate values for the solution of the initial-value problem $y'=f(x,y)$, $y(x_0)=y_0$, with step size $h$, at $x_n=x_{n-1}+h$, are the solutions of the equations
\[y_n=y_{n-1}+hf(x_n,y_n),\quad n\in\bbN.\]
\end{document}
